\documentclass[12pt, a4paper]{article}

\usepackage[utf8]{inputenc}
\usepackage[T2A]{fontenc}
\usepackage[russian]{babel}
\usepackage{amsmath, amssymb, amsfonts}
\usepackage{booktabs}
\usepackage{tabularx}
\usepackage{multirow}
\usepackage{threeparttable}
\usepackage{geometry}
\geometry{left=30mm, right=15mm, top=20mm, bottom=20mm}
\usepackage{setspace}
\onehalfspacing
\setlength{\emergencystretch}{3em}
\usepackage{caption}
\usepackage{float}
\usepackage{graphicx}
\usepackage{xcolor}
\usepackage{framed}
\usepackage{indentfirst}
\usepackage{enumitem}
\usepackage{cite}

\graphicspath{{Analysis/assets/}{../Analysis/assets/}}

\captionsetup{font=small, labelsep=period}
\captionsetup[figure]{name=Рисунок}

\setlist[itemize]{itemsep=0pt, topsep=0pt, parsep=0pt, partopsep=0pt}
\setlist[enumerate]{itemsep=0pt, topsep=0pt, parsep=0pt, partopsep=0pt}

\renewenvironment{quote}{\begin{leftbar}\noindent\ignorespaces}{\end{leftbar}}

% Два рисунка в строку, у каждого своя подпись и источник
\newcommand{\figurepair}[6]{
\begin{figure}[H]
    \centering
    \begin{minipage}[t]{0.49\textwidth}
        \centering
        \includegraphics[width=\linewidth, height=5.4cm, keepaspectratio]{#1}
        \caption{#2}
        \label{#3}
        \vspace{-2mm}
        {\footnotesize Источник: составлено авторами по данным X5 Group, Росстата и GeoNames.}
    \end{minipage}
    \hfill
    \begin{minipage}[t]{0.49\textwidth}
        \centering
        \includegraphics[width=\linewidth, height=5.4cm, keepaspectratio]{#4}
        \caption{#5}
        \label{#6}
        \vspace{-2mm}
        {\footnotesize Источник: составлено авторами по данным X5 Group, Росстата и GeoNames.}
    \end{minipage}
\end{figure}
}

% Широкий рисунок на всю строку
\newcommand{\figurewide}[3]{
\begin{figure}[H]
    \centering
    \includegraphics[width=\textwidth]{#1}
    \caption{#2}
    \label{#3}
    \vspace{-2mm}
    {\footnotesize Источник: составлено авторами по данным X5 Group, Росстата и GeoNames.}
\end{figure}
}

% Отдельный крупный рисунок для ключевых карт
\newcommand{\figuresingle}[3]{
\begin{figure}[H]
    \centering
    \includegraphics[width=0.75\textwidth]{#1}
    \caption{#2}
    \label{#3}
    \vspace{-2mm}
    {\footnotesize Источник: составлено авторами по данным X5 Group, Росстата и GeoNames.}
\end{figure}
}

% Рисунок раздела о прогнозе: подпись, метка и свой источник
\newcommand{\figureforecast}[5]{
\begin{figure}[H]
    \centering
    \includegraphics[width=#2\textwidth]{#1}
    \caption{#3}
    \label{#4}
    \vspace{-2mm}
    {\footnotesize Источник: #5.}
\end{figure}
}

\newcommand{\tablesourcetext}[1]{\par\vspace{1mm}{\footnotesize Источник: #1.}}

\newcommand{\tablesource}{\par\vspace{1mm}{\footnotesize Источник: расчеты авторов по данным X5 Group, Росстата и GeoNames.}}

\begin{document}

\begin{titlepage}
    \thispagestyle{empty}
    \begin{center}
        \linespread{1.1}\selectfont
        \textbf{Домашнее задание №1}\\[1cm]
        \Large{\textbf{Пространственная эконометрика}}\\[1cm]
        \large{Пространственные эффекты в выручке продуктового ритейла:\\
        моделирование розничного товарооборота магазинов сети\\
        «Пятерочка» по локальным рынкам России}\\[1.5cm]
        \normalsize
        \textbf{Выполнили:}\\[0.4cm]
        Гончаров Алексей\\
        Дугарев Денис\\
        Русских Никита\\
        Смирнов Владимир\\
        Хмельницкий Богдан\\[2cm]
    \end{center}
    \vfill
    \begin{center}
        Москва 2026
    \end{center}
\end{titlepage}

\setcounter{page}{2}

\section*{Структура работы}

Работа состоит из двух частей:

\begin{enumerate}
    \item Часть I - исследование пространственных эффектов в выручке магазинов сети «Пятерочка» на
    локальных рынках России методами пространственной эконометрики (с. \pageref{sec:research});
    \item Часть II - прогноз розничного товарооборота магазинов сети на март 2025 года, построенный
    по итогам соревнования X5 Group и скорректированный с использованием методов пространственной
    эконометрики (с. \pageref{sec:forecast}).
\end{enumerate}

\clearpage

\part*{Часть I. Исследование пространственных эффектов в выручке}
\label{sec:research}

% ============================================================
\section*{Введение}

Выручка продуктового магазина определяется не только его собственными характеристиками, но и
тем, что происходит вокруг. Близость крупного города, оживленных трасс и развитой торговой среды
расширяет поток покупателей и повышает выручку. Соседние магазины, в том числе принадлежащие той
же сети, делят этот поток между собой и выручку снижают. Оба эффекта возможны одновременно, и
заранее неизвестно, какой из них преобладает.

Цель исследования состоит в том, чтобы объяснить различия в выручке магазинов сети «Пятерочка»
между локальными рынками России (населенными пунктами, в которых сеть представлена несколькими
магазинами) с учетом взаимного влияния соседних рынков и выяснить, каким механизмом это влияние
объясняется. Поставлены следующие задачи:

\begin{enumerate}
    \item Собрать данные о выручке и характеристиках магазинов на уровне населенных пунктов;
    \item Установить, образует ли выручка пространственные кластеры;
    \item Оценить связь выручки населенного пункта с характеристиками и выручкой соседних
    населенных пунктов и разделить ее на прямые и косвенные эффекты;
    \item Выяснить, сохраняется ли связь с соседями при учете региональных различий и во времени;
    \item Проверить, различаются ли эффекты по территории.
\end{enumerate}

Актуальность работы связана с тем, что продуктовая розница является одной из крупнейших отраслей
российской экономики, а крупные сети принимают решения об открытии магазинов с опорой на прогноз
выручки. Если выручка точки зависит от окружения на десятки километров вокруг, прогноз и выбор
места без учета пространственных эффектов оказываются смещенными.

\clearpage

% ============================================================
\section*{Обзор литературы}

Пространственная организация розничной торговли изучается в экономике отраслевых рынков. Д.Л. Хафф предложил
гравитационную модель зоны притяжения магазина, в которой вероятность выбора покупателем из места
$i$ магазина $j$ растет с привлекательностью магазина $S_j$ и убывает с расстоянием $d_{ij}$
\cite{huff1964}:
\[
P_{ij} = \frac{S_j^{\alpha} d_{ij}^{-\lambda}}{\sum_k S_k^{\alpha} d_{ik}^{-\lambda}}.
\]
\indent Ожидаемая выручка магазина равна сумме спроса окрестных мест, взвешенной этими
вероятностями. Из модели следуют три предсказания, которые проверяются в настоящей работе:

\begin{enumerate}
    \item Выручка растет с численностью населения и потоком покупателей вокруг магазина;
    \item При прочих равных появление нового конкурента поблизости снижает вероятность выбора
    рассматриваемого магазина и его ожидаемую выручку;
    \item Спрос соседних населенных пунктов перетекает в магазин с весом, убывающим с расстоянием,
    что прямо соответствует матрице пространственных весов с затуханием по расстоянию и
    пространственным лагам характеристик соседей.
\end{enumerate}

Т.Дж. Холмс исследовал распространение сети Walmart в США и показал, что сеть открывала новые
магазины вплотную к существующим, поскольку плотная сеть снижает логистические издержки
распределительных центров, и сознательно принимала каннибализацию выручки собственных магазинов
\cite{holmes2011}. Для настоящей работы из этого следуют две проверки:

\begin{enumerate}
    \item Если каннибализация существенна, плотность собственных магазинов сети рядом с точкой
    должна быть связана с меньшей выручкой;
    \item Если же преобладает выгода плотности, плотность сети в соседних населенных пунктах может
    быть связана с выручкой положительно.
\end{enumerate}

П. Цзя оценила модель входа на локальные рынки сетей Kmart и Walmart и небольших дискаунтеров по
округам США и показала, что вход крупной сети вытесняет небольших конкурентов, а решения сети на
соседних рынках взаимосвязаны через выгоды от плотности \cite{jia2008}. Эта работа обосновывает
рассмотрение локальных рынков как связанных между собой, а не независимых наблюдений, и указывает
на эндогенность размещения магазинов и конкурентов, которая ограничивает причинную интерпретацию
оценок. В перечисленных работах конкуренция и плотность сети изучаются на американских данных
средствами моделей входа и размещения, а выручка магазинов не моделируется
пространственно-эконометрическими методами. Настоящее исследование переносит этот вопрос на
российскую продуктовую розницу и оценивает, как выручка связана с конкуренцией, плотностью сети и
характеристиками соседних рынков.

Для российских данных пространственно-эконометрический инструментарий применили О.А. Демидова и
Е.А. Тимофеева при оценке кривой заработной платы по 81 региону за 2005-2018 годы
\cite{demidova2021}. Без учета взаимного влияния регионов эффект безработицы на заработную плату
оказался завышен примерно вдвое, а разложение предельных эффектов в моделях SAR и SDM показало,
что рост безработицы в регионе снижает заработную плату и в соседних регионах. Для настоящей
работы этот результат важен по двум причинам:

\begin{enumerate}
    \item Заработная плата и безработица в субъекте входят в модель выручки как показатели
    платежеспособного спроса и сами пространственно связаны;
    \item Пропуск пространственной зависимости может так же смещать оценки влияния
    характеристик рынка на выручку.
\end{enumerate}

\clearpage

% ============================================================
\section*{Данные}

Первичными данными выступает помесячная панель выручки 18 657 магазинов сети «Пятерочка» с
января 2023 по февраль 2025 года, предоставленная X5 Group участникам соревнования по
прогнозированию розничного товарооборота\footnote{Значения выручки могли быть преобразованы
организаторами соревнования при обезличивании. Пусть истинная выручка равна $R_{it}$, а
опубликованная $R^{*}_{it} = c R_{it}$, где $c > 0$ - общий множитель. Тогда
$\ln R^{*}_{it} = \ln R_{it} + \ln c$, и слагаемое $\ln c$ поглощается константой регрессии, не
меняя остальных коэффициентов.}. Для каждого магазина известны населенный пункт, регион и
социально-географические характеристики, в том числе численность населения и домохозяйств в зоне
охвата, пеший и автомобильный трафик, число конкурентов и объектов инфраструктуры в заданном
радиусе, число касс, категории торговой площади и возраста магазина. Первичные данные дополнены собранной
самостоятельно открытой статистикой и координатами.

Данные сети содержат несколько видов искажений, которые исправлялись при подготовке среза:

\begin{enumerate}
    \item Предполагается, что некоторые исходные показатели содержат ошибки измерения: например,
    численность населения различается между магазинами одного населенного пункта. Трафик и
    характеристики окружения фиксируются на уровне отдельных торговых точек и могут отражать
    как реальные различия, так и неточности измерения. Агрегирование снижает чувствительность к
    отдельным нетипичным значениям;
    \item Нулевые значения в признаках, которые у работающего магазина не могут быть нулевыми,
    обозначают пропуски и в расчет не включаются. Таких значений 1128 для пешего трафика,
    698 для автомобильного, 780 для численности населения, 143 для числа касс и 32 для числа
    домохозяйств;
    \item Пятнадцать населенных пунктов записаны в данных по-разному (например, «Малаховка пгт.»
    и «Малаховка рп»), и их написания объединены;
    \item Пять магазинов, у которых вместо населенного пункта указан район, исключены;
    \item У 203 магазинов выручка меняется между соседними месяцами более чем втрое. Эти магазины
    сохранены, а результаты без них приводятся в проверке устойчивости.
\end{enumerate}

Единицей наблюдения выбран населенный пункт (локальный рынок), что обосновано тремя соображениями:

\begin{enumerate}
    \item Координаты известны только на уровне населенного пункта. Присвоение этих координат
    отдельным магазинам создало бы совпадающие точки и нулевые расстояния внутри города, и
    матрица весов не отражала бы реального расположения магазинов. Расстояния между населенными
    пунктами при этом измеряются корректно.
    \item Агрегирование снижает шум. Медиана непрерывных признаков, среднее счетных показателей
    окружения и доли категорий дают устойчивое описание торговой среды населенного пункта.
    Снижается прежде всего индивидуальный шум отдельных магазинов, а общая для всех магазинов
    ошибка измерения агрегированием не устраняется.
    \item Этот уровень соответствует исследовательскому вопросу. Для большинства покупателей
    рынок повседневных продуктов ограничен своим населенным пунктом и ближайшими соседними, и
    именно связь характеристик населенного пункта и его соседей с выручкой магазинов изучается в
    работе.
\end{enumerate}

В выборку включены населенные пункты, в которых сеть представлена не менее чем тремя магазинами.
При одном или двух магазинах характеристики населенного пункта определяются шумом и форматом
отдельной точки, а медиана трех значений уже отбрасывает одно крайнее. Порог исключает 2422
небольших населенных пункта с 2969 магазинами и оставляет 933 населенных пункта из 67 регионов
с 15 683 магазинами, то есть 84\% магазинов сети в данных. Среди них 725 городов, 120 поселков и
поселков городского типа, 88 сел, станиц и деревень. Отбор зависит от решений сети о
размещении, поэтому выводы относятся к рынкам, на которых сеть уже присутствует несколькими
магазинами, и не переносятся на все населенные пункты России. Чувствительность к порогу
проверяется на подвыборке с пятью и более магазинами. Состав магазинов внутри населенного пункта
за период не меняется, и панель строго сбалансирована, $N = 933$, $T = 26$.

Численность населения сети зашумлена и содержит нули, поэтому она заменена данными бюллетеня
Росстата по муниципальным образованиям на 1 января 2024 года для 822 населенных пунктов. Бюллетень
охватывает только города и поселки городского типа. Для 111 сел и станиц, для которых отсутствуют
сведения в бюллетене, численность населения определена как медиана значений, указанных в исходных
данных сети для магазинов соответствующего населенного пункта. Мы допускаем такую замену, поскольку
там, где известны оба источника, медиана данных сети отличается от данных Росстата в среднем на
1.5\%.

Росстат публикует уровень безработицы по субъектам только за скользящие трехмесячные периоды,
поскольку выборка обследования рабочей силы в отдельном субъекте за один месяц мала. Каждый период
отнесен к своему центральному месяцу, например период «январь - март 2023 года» дает значение за
февраль 2023 года. Так получается помесячный ряд, среднее которого за 2024 год совпадает с
годовым значением с точностью до 0.3 процентного пункта.

Выручка переведена в цены января 2023 года региональным индексом цен на продовольственные товары,
поскольку основную часть ассортимента продуктового магазина составляет продовольствие:
\[
RTO^{real}_{it} = \frac{RTO_{it}}{P^{food}_{r(i),t}}, \qquad
P^{food}_{r,t} = \prod_{s = \text{фев 2023}}^{t} \frac{\pi^{food}_{r,s}}{100},
\]
где $\pi^{food}_{r,s}$ - индекс цен на продовольствие в субъекте $r$ в месяце $s$ в процентах к
предыдущему месяцу. Реальная заработная плата получена делением на общий индекс цен. К февралю
2025 года продовольственный индекс вырос на 15-27\% в зависимости от субъекта. Цепной индекс
устраняет инфляцию во времени, но не различия в уровне цен между регионами на базовую дату,
поэтому сопоставление выручки разных субъектов дополнительно проверяется спецификацией с
фиксированными эффектами субъектов.

Зависимая переменная - средний за 26 месяцев логарифм реальной выручки магазина в населенном
пункте:
\[
y_i = \frac{1}{26} \sum_{t=1}^{26} \ln \left( \frac{1}{n_i} \sum_{k \in i} RTO^{real}_{kt} \right),
\]
где $n_i$ - число магазинов сети в населенном пункте $i$. Экономически $y_i$ описывает типичный
многолетний уровень месячной выручки магазина в населенном пункте. Среднее логарифмов равно
логарифму геометрического среднего выручки, поэтому отдельные месяцы с аномально высокой или
низкой выручкой влияют на него слабее, чем на арифметическое среднее. Усреднение по двум годам
сглаживает месячные шоки и сезонность, поскольку почти каждый календарный месяц входит в окно
одинаковое число раз. Исключение составляют январь и февраль, которые входят в окно три раза, а не
два, и сезонность из-за этого устраняется не полностью. Поэтому в проверке устойчивости
используется окно ровно из 24 месяцев, с марта 2023 по февраль 2025 года. Описательные
статистики показателей приведены в таблице \ref{tab:descriptive}.

\begin{table}[H]
    \centering
    \footnotesize
    \renewcommand{\arraystretch}{1.1}
    \caption{Показатели исследования и описательные статистики, $N = 933$}
    \label{tab:descriptive}
    \begin{tabularx}{\textwidth}{>{\raggedright\arraybackslash}X >{\raggedright\arraybackslash}p{4.3cm} rrrr}
        \toprule
        \textbf{Показатель} & \textbf{Единица измерения} & \textbf{Среднее} & \textbf{Ст. откл.} & \textbf{Мин.} & \textbf{Макс.} \\
        \midrule
        $y$, выручка магазина & средний логарифм рублей в месяц & 18.077 & 0.295 & 17.393 & 19.082 \\
        Численность населения & логарифм человек & 10.397 & 1.198 & 6.498 & 16.392 \\
        Домохозяйства в зоне охвата & логарифм единиц & 7.427 & 0.814 & 4.905 & 9.450 \\
        Пеший трафик & логарифм человек в час & 4.724 & 0.360 & 3.597 & 6.005 \\
        Автомобильный трафик & логарифм автомобилей в час & 4.972 & 0.616 & 2.674 & 7.533 \\
        Продуктовые конкуренты & магазинов в радиусе 500 м & 3.320 & 1.605 & 0 & 11.000 \\
        Плотность сети & магазинов сети в радиусе 500 м & 0.631 & 0.551 & 0 & 3.667 \\
        Пункты выдачи маркетплейсов & в радиусе 100 м & 1.645 & 0.825 & 0 & 6.000 \\
        Медицинские учреждения и аптеки & в радиусе 300 м & 1.338 & 1.008 & 0 & 7.000 \\
        Кассы & штук на магазин & 7.333 & 1.492 & 3.333 & 14.750 \\
        Доля больших магазинов & доля & 0.191 & 0.193 & 0 & 1 \\
        Доля новых магазинов & доля & 0.212 & 0.212 & 0 & 1 \\
        Реальная зарплата в субъекте & логарифм рублей & 11.013 & 0.209 & 10.604 & 11.889 \\
        Безработица в субъекте & процентов & 2.472 & 0.745 & 1.381 & 9.212 \\
        \bottomrule
    \end{tabularx}
    \tablesource
\end{table}

Объясняющие переменные частично дублируют друг друга, например численность населения и число
домохозяйств в зоне охвата связаны между собой. Степень такого дублирования измеряется
коэффициентом увеличения дисперсии (VIF), который показывает, во сколько раз дисперсия оценки
коэффициента в линейной регрессии выше, чем при отсутствии связи регрессора с остальными.
Наибольшее значение равно 3.11 (число домохозяйств), что указывает на умеренную
мультиколлинеарность.

Сильнее всего с выручкой связаны число касс ($r = 0.71$), реальная зарплата в субъекте
($r = 0.50$) и доля больших магазинов ($r = 0.47$). Принадлежность к субъекту объясняет 58\% дисперсии зависимой
переменной, то есть регрессия выручки только на индикаторы субъектов дает $R^2 = 0.58$, и большая
часть различий между населенными пунктами приходится на различия средних уровней регионов. Это
означает, что соседние населенные пункты могут быть похожи из-за общей региональной среды, а не
из-за влияния друг на друга, поэтому в работе проводятся проверки с учетом региональных эффектов.

Три показателя чека, содержащиеся в данных (среднее число товаров, промо-товаров и отмен), в
модель не включаются. Они отражают сами покупки, из которых складывается выручка, и определяются
одновременно с ней как результат спроса, а не как его внешняя причина. Их включение сделало бы
эти переменные эндогенными и затруднило бы интерпретацию остальных эффектов.

\subsection*{Предварительный анализ данных}

Населенные пункты выборки сосредоточены в европейской части России, на Урале и в южной Сибири
(рисунок \ref{fig:location}). Крупнейшие узлы сети - Москва с 1278 магазинами и Санкт-Петербург с
498 магазинами, медиана составляет 6 магазинов на населенный пункт.

\figurewide{settlements_russia.png}{Населенные пункты выборки, размер точки - число магазинов сети}{fig:location}

На карте выручки высокие значения образуют сплошную зону вокруг Москвы и Санкт-Петербурга, а
низкие - на юге и в Поволжье (рисунок \ref{fig:revenue_map}). Судя по гистограмме, распределение
зависимой переменной в логарифмической шкале не имеет выраженной асимметрии (рисунок
\ref{fig:distribution}). Средняя реальная выручка магазина по населенным пунктам составляет от 36
до 194 млн рублей в месяц при медиане 69 млн.

\figurepair{revenue_map.png}{Средний месячный логарифм реальной выручки магазина}{fig:revenue_map}
{revenue_distribution.png}{Распределение логарифма реальной выручки магазина по населенным пунктам}{fig:distribution}

Номинальная выручка имеет выраженную сезонность с пиком в декабре и минимумом в начале года
(рисунок \ref{fig:dynamics}; средняя выручка магазина рассчитана по населенным пунктам с равными весами, а
не по всем магазинам сети). В реальном выражении январь 2024 года выше января 2023 года на 5\%,
январь 2025 года выше января 2024 года на 2\%, а февраль 2025 года ниже февраля 2024 года на 4\%.
Во всех месяцах наблюдается положительная пространственная автокорреляция сопоставимой величины.
Индекс Морана составляет 0.51-0.61, а оценки пространственного коэффициента месячных SDM -
0.45-0.53 (рисунок \ref{fig:monthly}). Это согласуется с использованием усредненного за период
показателя для анализа различий между населенными пунктами, однако не доказывает неизменность
расположения кластеров, поскольку при похожем значении индекса состав кластеров может меняться.
Логарифмирование выручки обосновано отдельно. Оно переводит различия между населенными пунктами
в относительную шкалу, уменьшает влияние крупнейших рынков и позволяет интерпретировать
коэффициенты при логарифмированных регрессорах как эластичности.

\figurepair{revenue_dynamics.png}{Средняя выручка магазина по месяцам}{fig:dynamics}
{moran_monthly.png}{Индекс Морана выручки и коэффициент $\rho$ SDM, оцененные по каждому месяцу}{fig:monthly}

\clearpage

% ============================================================
\section*{Модель}

\subsection*{Обоснование набора объясняющих переменных}

Набор регрессоров следует из модели зоны притяжения магазина \cite{huff1964}. Выручка растет с
размером рынка и потоком покупателей, снижается с числом конкурентов и зависит от формата точки
и доходов населения. Размер рынка измеряется численностью населения населенного пункта и числом
домохозяйств в зоне охвата магазина. Поток покупателей измеряется пешим и автомобильным трафиком.
Конкуренция описывается числом продуктовых магазинов и магазинов самой сети в радиусе 500 м, а
также числом пунктов выдачи маркетплейсов, через которые часть спроса может уходить в электронную
торговлю. Окружение описывается числом медицинских учреждений и аптек как показателем деловой
активности, платежеспособный спрос - реальной заработной платой и уровнем безработицы в субъекте.

Формат магазина задается числом касс и долями больших и новых магазинов. Выручка на магазин
механически растет с размером точки, поэтому без учета формата сравнение населенных пунктов смешивало бы
пространственные эффекты со структурой форматов сети. Вместе с тем число касс и торговую площадь
сеть выбирает с учетом ожидаемого спроса, и эти переменные частично эндогенны. Их включение
меняет объект оценки, поскольку основная спецификация сравнивает выручку магазинов одного
формата, а направление смещения из-за эндогенности контролей заранее неизвестно. Поэтому
спецификации с контролями формата и без них рассматриваются как проверка чувствительности к
набору контролей.

Пространственные лаги $WX$ строятся для шести показателей спроса и конкуренции, которые
содержательно могут переливаться между соседними рынками. Это численность населения, число
домохозяйств, пеший и автомобильный трафик, число конкурентов и плотность сети. Лаги зарплаты и
безработицы не добавляются, поскольку у соседей из того же субъекта эти показатели почти
совпадают со своими.

\subsection*{Матрицы пространственных весов}

Матрица весов $W = (w_{ij})$ задает, какие населенные пункты считаются соседями и насколько сильно
они связаны. Во всех вариантах $w_{ii} = 0$, $w_{ij} \ge 0$, и матрица нормирована по строкам,
$w_{ij} = \tilde{w}_{ij} / \sum_k \tilde{w}_{ik}$, так что $Wy$ равен средневзвешенному значению
показателя у соседей. Расстояния $d_{ij}$ вычисляются по большому кругу между координатами
населенных пунктов. Рассматриваются шесть вариантов исходных весов $\tilde{w}_{ij}$:

\begin{enumerate}
    \item $\tilde{w}_{ij} = \exp(-d_{ij} / \delta)$ при $d_{ij} < 300$ км, $\delta = 50$ км
    (экспоненциальное затухание).
    \item $\tilde{w}_{ij} = 1$, если $j$ входит в число 8 ближайших соседей $i$ или $i$ входит в
    число 8 ближайших соседей $j$ (симметризованная матрица $k$ ближайших соседей).
    \item $\tilde{w}_{ij} = d_{ij}^{-1}$ при $d_{ij} < 200$ км (обратные расстояния).
    \item $\tilde{w}_{ij} = d_{ij}^{-2}$ при $d_{ij} < 400$ км (обратные квадраты расстояний).
    \item $\tilde{w}_{ij} = 1$ при $d_{ij} < 100$ км (пороговая матрица).
    \item $\tilde{w}_{ij} = 1$, если $j$ входит в число 10 населенных пунктов, ближайших к $i$ по логарифму
    численности населения (экономическая близость без учета географии).
\end{enumerate}

Во всех матрицах, кроме пороговой, населенный пункт без соседей в радиусе отсечения получает ближайшего
соседа. Пороговая матрица оставлена в исходном виде, чтобы показать последствия изолированных
населенных пунктов. Экономическая матрица служит проверкой чувствительности к одной альтернативной структуре
связей - сходству населенных пунктов по численности населения. Отсутствие зависимости на ней не исключает
общих региональных факторов, транспортной доступности и управления сетью как причин
географического паттерна.

\textbf{Матрицы пространственных весов} сравниваются по максимизированному
логарифмическому правдоподобию отдельно в моделях SAR и SDM \cite{elhorst2014, lesage2009}.
При смене матрицы внутри каждой модели сохраняются выборка, набор регрессоров и число
оцениваемых параметров. Поэтому большее значение критерия указывает на лучшее соответствие
данной спецификации наблюдаемым данным. Матрица с экспоненциальным затуханием весов
$\tilde{w}_{ij} = \exp(-d_{ij}/\delta)$ ($\delta = 50$ км) при $d_{ij} < 300$ км и нулевыми
весами за пределами этого расстояния обеспечивает наибольшее правдоподобие в обеих моделях и
выбирается основной (таблица \ref{tab:weights}). Масштаб затухания $\delta$ дополнительно
проверяется по профилю правдоподобия SDM на сетке значений. Поскольку выбор производится по
исследуемым данным, он не устанавливает истинную структуру пространственных связей или границу
локального рынка. Построчная нормировка делает веса относительными - у населенного пункта с единственным
соседом в 150 км суммарный вес соседей такой же, как у населенного пункта в плотной агломерации. Это
ограничение учитывается при интерпретации.

\subsection*{Статистики пространственной автокорреляции}

Глобальная автокорреляция измеряется индексом Морана, индексом Гири и статистикой Гетиса-Орда:
\[
I = \frac{n}{\sum_{i,j} w_{ij}} \cdot \frac{\sum_{i,j} w_{ij} (y_i - \bar{y})(y_j - \bar{y})}{\sum_i (y_i - \bar{y})^2}, \qquad
C = \frac{(n - 1) \sum_{i,j} w_{ij} (y_i - y_j)^2}{2 \sum_{i,j} w_{ij} \sum_i (y_i - \bar{y})^2}, \qquad
G = \frac{\sum_{i \ne j} w_{ij} x_i x_j}{\sum_{i \ne j} x_i x_j}.
\]
\indent Статистика $G$ определена для положительных уровней показателя и бинарной матрицы
\cite{getis1992}, поэтому она считается по реальной выручке $x_i$ в рублях и соседству в радиусе
100 км. Двумерный индекс Морана $I_{xy} = \sum_{i,j} w_{ij} z^x_i z^y_j / \sum_i (z^x_i)^2$, где
$z$ - стандартизованные значения, измеряет связь показателя $x$ в населенном пункте с показателем $y$ у
соседей. Локальные статистики разлагают глобальные на вклады отдельных населенных пунктов:
\[
I_i = z_i \sum_j w_{ij} z_j, \qquad
C_i = \sum_j w_{ij} (z_i - z_j)^2, \qquad
G^*_i = \frac{\sum_j w^*_{ij} y_j}{\sum_j y_j},
\]
где в $G^*_i$ матрица $w^*_{ij}$ совпадает с основной, но включает сам населенный пункт с весом
$\exp(0) = 1$. Многомерный локальный индекс Гири суммирует $C_i$ по нескольким стандартизованным
показателям и проверяет сходство населенного пункта с соседями в пространстве этих признаков, а не по каждому
из них отдельно. Значимость локальных статистик определяется перестановочным тестом с 9999
перестановками и двусторонней альтернативой. Его $p$-значение равно удвоенной доле перестановок в
меньшем хвосте распределения, $p_i = 2 (\min(m^-_i, m^+_i) + 1) / (R + 1)$. Поправка на
множественное тестирование выполняется методом Бенджамини-Хохберга \cite{benjamini1995},
контролирующим ожидаемую долю ложноположительных результатов среди всех результатов,
признанных статистически значимыми (FDR). Метод гарантирует этот контроль при независимых или
положительно зависимых статистиках, а локальные статистики соседних населенных пунктов зависимы, поэтому для
локального индекса Морана дополнительно приводится поправка Бенджамини-Иекутиели
\cite{benjamini2001}, верная при любой зависимости. Поправка Бонферрони при 933 тестах требует
$p < 0.00005$, что при таком числе перестановок недостижимо, поэтому она сравнивается с FDR на
аналитических $p$-значениях.

\subsection*{Модели для данных типа cross-section и выбор спецификации}

Спецификация выбирается двумя подходами, которые дополняют друг друга. Подход от простой модели
к сложной начинается с МНК и добавляет пространственную компоненту, на которую указывают тесты
по остаткам \cite{anselin1996}. Подход от общей модели к частной начинается с модели, содержащей
все пространственные взаимодействия, и последовательно проверяет ограничения
\cite{elhorst2014}. Если подходы приводят к разным моделям, выбор между ними делается тестом
отношения правдоподобия, когда одна модель вложена в другую, и по информационным критериям и
остаточной автокорреляции в остальных случаях.

В подходе от простой модели к сложной оценивается МНК, и по его остаткам вычисляются тесты
множителей Лагранжа $LM_{lag}$ ($H_0: \rho = 0$ против $H_1: \rho \ne 0$) и $LM_{err}$
($H_0: \lambda = 0$ против $H_1: \lambda \ne 0$). Робастные версии тестов устойчивы к присутствию
конкурирующей пространственной альтернативы, но не к гетероскедастичности. Если значим только один
из робастных тестов, выбирается соответствующая модель SAR или SEM.

В подходе от общей модели к частной исходной служит модель GNS, содержащая все три вида
пространственных взаимодействий:
\[
y = \rho W y + \alpha \iota_n + X \beta + W X^{D} \theta + u, \qquad u = \lambda W u + \varepsilon,
\]
где $X^{D}$ - шесть регрессоров с пространственными лагами. Остальные модели получаются из нее
ограничениями на параметры:

\begin{enumerate}
    \item Модель Дарбина SDM, $\lambda = 0$.
    \item Модель Дарбина с пространственной ошибкой SDEM, $\rho = 0$.
    \item Модель SARAR с лагом зависимой переменной и пространственной ошибкой, $\theta = 0$.
    \item Модель пространственного лага SAR, $\theta = 0$ и $\lambda = 0$.
    \item Модель с лагами регрессоров SLX, $\rho = 0$ и $\lambda = 0$.
    \item Модель пространственной ошибки SEM, вложенная в SDM с лагами всех 13 регрессоров через
    ограничение общего множителя $\theta = -\rho \beta$ \cite{elhorst2014}.
    \item МНК, $\rho = \theta = \lambda = 0$.
\end{enumerate}

Ограничения проверяются тестом отношения правдоподобия
\[
LR = 2 \left( \ln L_{\text{общая}} - \ln L_{\text{ограниченная}} \right) \sim \chi^2_q,
\]
где $q$ - число ограничений, а $H_0$ состоит в выполнении проверяемых ограничений. Параметры
$\rho$ и $\lambda$ в GNS с одной матрицей весов идентифицируются слабо, поэтому GNS используется
только как отправная точка для тестов. Невложенные модели сравниваются по информационным
критериям AIC и BIC. Остаточная автокорреляция проверяется тестом Морана для МНК и SLX и
LM-тестом для моделей с лагом зависимой переменной, гетероскедастичность - тестом Бройша-Пагана
($H_0$ - гомоскедастичность ошибок).

Модели оцениваются методом максимального правдоподобия. Пространственный лаг $Wy$ коррелирует с
ошибкой, поскольку выручка соседей сама зависит от выручки рассматриваемого населенного пункта и,
следовательно, от его случайного шока, и МНК для таких моделей несостоятелен. Поэтому оценки SAR и
SDM дополнительно проверяются пространственным двухшаговым МНК с инструментами $X$, $WX$, $W^2 X$
\cite{kelejian1998} и стандартными ошибками Уайта. Этот метод инструментирует только $Wy$ и
предполагает экзогенность $X$, и если число конкурентов или касс эндогенно, $WX$ и $W^2 X$
наследуют эту проблему. Оценка SEM дополнительно проверяется обобщенным МНК Келеджиана-Пручи
\cite{kelejian1999}, в котором данные преобразуются умножением на $(I - \hat{\lambda} W)$, и его
версией, робастной к гетероскедастичности \cite{kelejian2010}. LR-тесты построены при однородной
дисперсии ошибок, поэтому ключевое ограничение $\theta = 0$ дополнительно проверяется тестом
Вальда с ковариационной матрицей двухшагового МНК, робастной к гетероскедастичности.

\subsection*{Прямые, косвенные и общие эффекты выбранной модели}

В моделях с лагом зависимой переменной коэффициенты не являются предельными эффектами, поэтому
выбранная модель интерпретируется через разложение эффектов \cite{lesage2009}. В SDM предельные
эффекты регрессора $x_m$ задаются матрицей
\[
S_m(W) = \frac{\partial E(y)}{\partial x_m'} = (I - \rho W)^{-1} (\beta_m I + \theta_m W).
\]
\indent Средний прямой эффект равен среднему диагональному элементу, общий - средней сумме строки,
косвенный - их разнице:
\[
ADE_m = \frac{1}{n} \operatorname{tr} S_m, \qquad
ATE_m = \frac{1}{n} \iota' S_m \iota, \qquad
AIE_m = ATE_m - ADE_m.
\]
\indent Косвенные эффекты SAR и SDM являются глобальными спилловерами и распространяются на
соседей всех порядков через разложение $S_m = \sum_{q \ge 0} (\rho W)^q (\beta_m I + \theta_m W)$.
В SLX и SDEM косвенный эффект равен $\theta_m$ и является локальным спилловером, затрагивающим
только прямых соседей \cite{vega2015}. Эффекты выбранной модели сравниваются с эффектами SAR, SLX
и SDEM. Значимость эффектов оценивается по 2000 симуляциям из асимптотического распределения
оценок параметров.

\subsection*{Дополнительные модели}

Модели для среза дополняются тремя группами моделей, каждая из которых отвечает на отдельный
вопрос:

\begin{enumerate}
    \item Пространственная панель проверяет, связаны ли месячные отклонения выручки соседних
    населенных пунктов после учета их постоянных различий;
    \item Динамическая панель оценивает инерционность выручки и то, передаются ли шоки от
    населенного пункта к населенному пункту во времени;
    \item Географически взвешенная регрессия проверяет, одинаковы ли эффекты по территории.
\end{enumerate}

Устойчивость пространственной структуры во времени
дополнительно проверяется оценкой SDM отдельно по каждому из 26 месяцев.

\subsubsection*{Пространственная панель}

Пространственная панель охватывает все 26 месяцев, $N = 933$, $T = 26$, 24 258 наблюдений. В ней
оцениваются модели с двунаправленными фиксированными эффектами
\[
y_{it} = \rho \sum_j w_{ij} y_{jt} + x_{it}' \beta + \mu_i + \tau_t + u_{it}, \qquad
u_{it} = \lambda \sum_j w_{ij} u_{jt} + \varepsilon_{it},
\]
где $x_{it}$ - реальная зарплата и безработица в субъекте. Индивидуальные эффекты $\mu_i$
поглощают все постоянные характеристики населенного пункта, включая формат магазинов, региональный уровень
цен и численность населения, а эффекты времени $\tau_t$ - общую сезонность и тренд. Коэффициенты
$\rho$ и $\lambda$ поэтому идентифицируются только по отклонениям выручки населенных пунктов от их
собственных средних, и межрегиональные различия уровней на них не влияют
\cite{lee2010, elhorst2014}. Выбор между фиксированными и случайными эффектами проверяется тестом
Хаусмана ($H_0$ - состоятельность оценок со случайными эффектами), выбор между лагом и ошибкой -
панельными LM-тестами и их робастными версиями, а также байесовскими апостериорными
вероятностями моделей \cite{lesage2014}. Зависимость между населенными пунктами проверяется CD-тестом
Песарана \cite{pesaran2021} ($H_0$ - отсутствие зависимости). Эффекты времени общие для страны,
поэтому дополнительно оценивается модель с эффектами «субъект $\times$ месяц», которые поглощают
и региональные месячные шоки.

\subsubsection*{Динамическая панель}

Динамическая модель включает лаг зависимой переменной:
\[
y_{it} = \gamma y_{i,t-1} + x_{it}' \beta + \mu_i + \tau_t + \varepsilon_{it}.
\]
\indent Лаг $y_{i,t-1}$ коррелирует с индивидуальным эффектом $\mu_i$. МНК без учета эффектов
завышает $\gamma$, а оценка с фиксированными эффектами ее занижает на величину порядка $1/T$
(смещение Никелла \cite{nickell1981}). Метод Ареллано-Бонда \cite{arellano1991} оценивает модель в
первых разностях обобщенным методом моментов с моментными условиями
\[
E\left( y_{i,t-s} \, \Delta \varepsilon_{it} \right) = 0, \qquad s \ge 2,
\]
а системный GMM \cite{blundell1998} добавляет уравнение в уровнях с инструментами
$\Delta y_{i,t-1}$. Чтобы число инструментов не росло с $T$, они сворачиваются
\cite{roodman2009}. Зарплата и безработица служат инструментами для самих себя, то есть считаются
строго экзогенными. Валидность инструментов проверяется $J$-тестом Хансена-Саргана ($H_0$ -
моментные условия выполняются), число степеней свободы которого равно разности числа инструментов
и числа параметров. Отсутствие автокорреляции первых разностей ошибок второго порядка
проверяется тестом AR(2) ($H_0$ - автокорреляции второго порядка нет). Долгосрочный эффект
регрессора в динамической модели равен $\beta / (1 - \gamma)$. Пространственная динамическая модель
\[
y_{t} = \gamma y_{t-1} + \rho W y_{t} + \eta W y_{t-1} + X_t \beta + W X_t \theta + \mu + \tau_t + \varepsilon_t
\]
оценивается методом максимального правдоподобия с двунаправленными фиксированными эффектами
\cite{yu2008, lee2010}. При конечном $T$ эта оценка смещена на величину порядка $1/T$, поэтому
применяется аналитическая коррекция смещения \cite{lee2010dynamic}. При ограничении
$\eta = -\gamma \rho$ переходная матрица $(I - \rho W)^{-1} (\gamma I + \eta W)$ равна $\gamma I$,
и модель SAR без $WX$ принимает вид
\[
y_t = \gamma y_{t-1} + (I - \rho W)^{-1} \left( X_t \beta + \mu + \tau_t \iota_n + \varepsilon_t \right).
\]
\indent При этом ограничении исчезает пространственная передача через прошлые значения, но
пространственный мультипликатор для текущих регрессоров и шоков сохраняется.

\subsubsection*{Географически взвешенная регрессия}

Пространственные модели предполагают одинаковые коэффициенты для всех населенных пунктов. Географически
взвешенная регрессия (GWR) \cite{fotheringham2002} ослабляет это предположение и оценивает для
каждого населенного пункта $i$ свою регрессию с весами, убывающими с расстоянием:
\[
y_i = \sum_k \beta_k(u_i, v_i) x_{ik} + \varepsilon_i, \qquad
\hat{\beta}(u_i, v_i) = (X' D_i X)^{-1} X' D_i y, \qquad
D_i = \operatorname{diag}\left( \exp\left( -\tfrac{1}{2} (d_{ij} / b_i)^2 \right) \right),
\]
где $(u_i, v_i)$ - координаты населенного пункта, а ширина окна $b_i$ адаптивная и равна расстоянию до
фиксированного числа ближайших населенных пунктов, которое выбирается кросс-валидацией. Постоянство каждого
коэффициента по территории проверяется F-тестом \cite{leung2000}. Значимость локального
коэффициента проверяется $t$-тестом в каждом населенном пункте с поправкой Бонферрони $0.05 / n$. Эти тесты
предполагают независимые ошибки, и поправка на множественность не исправляет стандартные ошибки
при пространственно зависимых остатках, поэтому локальные оценки GWR рассматриваются как
предварительные результаты, требующие осторожной интерпретации.

\clearpage

% ============================================================
\section*{Гипотезы}

\textbf{Гипотеза 1.} Выручка магазинов в соседних населенных пунктах схожа, и населенные пункты с
высокой и низкой выручкой образуют территориальные кластеры.
\[
H_0: I = E(I), \qquad H_1: I > E(I).
\]

Основанием служит модель зоны притяжения Д.Л. Хаффа \cite{huff1964}. В магазин приходят и
покупатели из соседних населенных пунктов, причем тем реже, чем дальше они живут. Поэтому у
соседних рынков частично общий спрос, и выручка на них должна быть похожей.

\textbf{Гипотеза 2.} Выручка в населенном пункте связана с выручкой соседних населенных пунктов и
тогда, когда населенные пункты одинаковы по населению, трафику, конкуренции и доходам, то есть
сходство выручки соседей не сводится к сходству их характеристик.
\[
H_0: \rho = 0, \qquad H_1: \rho \ne 0.
\]

П. Цзя показала, что соседние локальные рынки сети связаны между собой через выгоды от плотности
магазинов \cite{jia2008}. Знак связи заранее не задан. Положительный $\rho$ означает, что выручка
соседних рынков меняется в одну сторону, отрицательный - что покупатели перераспределяются между
ними.

\textbf{Гипотеза 3.} Чем больше рядом с магазином продуктовых конкурентов и магазинов самой сети,
тем ниже его выручка.
\[
H_0: ADE_{comp} = 0, \qquad H_1: ADE_{comp} < 0; \qquad
H_0: ADE_{dens} = 0, \qquad H_1: ADE_{dens} < 0.
\]

В модели Д.Л. Хаффа каждый новый конкурент забирает часть спроса \cite{huff1964}. Т.Дж. Холмс
показал, что сеть Walmart открывала магазины плотно и мирилась с тем, что они отбирают выручку
друг у друга \cite{holmes2011}. Эффекты конкурентов и плотности сети проверяются по отдельности.

\textbf{Гипотеза 4.} Выручка в населенном пункте зависит от численности населения, трафика,
конкуренции и плотности сети в соседних населенных пунктах.
\[
H_0: AIE_m = 0, \qquad H_1: AIE_m \ne 0.
\]

Направление влияния заранее неизвестно. Крупный соседний центр может приводить покупателей за
счет маятниковой миграции, а может оттягивать их к себе. Насыщенная торговая среда и плотная сеть
вокруг могут делить поток покупателей, а могут давать выгоды плотности, описанные Т.Дж. Холмсом и
П. Цзя \cite{holmes2011, jia2008}. Эффект проверяется для каждого показателя отдельно.

\textbf{Гипотеза 5.} Пространственная модель описывает различия выручки между населенными
пунктами лучше, чем обычная регрессия.
\[
H_0: \rho = 0 \text{ и } \theta = 0, \qquad H_1: \rho \ne 0 \text{ или } \theta \ne 0.
\]

О.А. Демидова и Е.А. Тимофеева показали на данных российских регионов, что без учета
пространственной зависимости оценки смещаются. В их работе эффект безработицы на заработную плату
оказался завышен примерно вдвое \cite{demidova2021}. Ограничения проверяются совместно.

\textbf{Гипотеза 6.} Месячные колебания выручки соседних населенных пунктов связаны между собой и
после учета постоянных различий между населенными пунктами и общих для страны сезонных колебаний.
\[
H_0: \rho = 0, \qquad H_1: \rho \ne 0 \quad \text{(SAR FE)}; \qquad
H_0: \lambda = 0, \qquad H_1: \lambda \ne 0 \quad \text{(SEM FE)}.
\]

Если соседние рынки получают общие шоки спроса, их отклонения от собственных средних уровней
движутся вместе. Такая связь может проявляться через лаг выручки или через ошибку, поэтому
гипотеза проверяется отдельно в моделях SAR FE и SEM FE.

\clearpage

% ============================================================
\section*{Результаты}

\subsection*{Выбор матрицы весов}

Все географические матрицы дают близкие значения индекса Морана, но различаются по качеству
описания данных (таблица \ref{tab:weights}). Наибольшее правдоподобие SAR и SDM достигается на
матрице с экспоненциальным затуханием. Профиль правдоподобия по масштабу затухания имеет пологий
максимум на 50-60 км (рисунок \ref{fig:profile}). Наибольшее значение 433.3 достигается при
60 км, при 50 км оно равно 432.3. Выбрано круглое значение 50 км, поскольку различие
правдоподобий пренебрежимо мало. Выбор матрицы по данным не учитывается в стандартных ошибках
последующих оценок, поэтому устойчивость выводов проверяется на других матрицах. При 50 км вес
соседа уменьшается вдвое на каждые 35 км, что правдоподобно для пригородной зоны, хотя по
поездкам покупателей этот масштаб не оценивался. После построчной нормировки один и тот же
масштаб означает разный охват, и средневзвешенное расстояние до соседей $\sum_j w_{ij} d_{ij}$
имеет медиану 89.5 км и меняется от 16.2 км в плотных агломерациях до 251 км. Для удаленных
населенных пунктов Севера и Сибири, где ближайший сосед находится дальше 100 км, единственный сосед получает
полный вес, и модель без таких населенных пунктов приводится в проверке устойчивости.

\begin{table}[H]
    \centering
    \footnotesize
    \caption{Сравнение матриц пространственных весов}
    \label{tab:weights}
    \begin{tabular}{lrrrrrr}
        \toprule
        \textbf{Матрица} & \textbf{\shortstack[r]{Среднее число\\соседей}} & \textbf{\shortstack[r]{Индекс\\Морана}} & \textbf{$\rho$ SAR} & \textbf{$\ln L$ SAR} & \textbf{$\rho$ SDM} & \textbf{$\ln L$ SDM} \\
        \midrule
        $\exp(-d/50)$, 300 км & 85.6 & 0.582 & 0.540 & 417.1 & 0.474 & 432.3 \\
        kNN, $k = 8$ & 9.9 & 0.585 & 0.532 & 408.6 & 0.467 & 426.9 \\
        $1/d$, 200 км & 46.7 & 0.577 & 0.532 & 412.4 & 0.475 & 425.5 \\
        $1/d^2$, 400 км & 131.7 & 0.592 & 0.494 & 408.2 & 0.461 & 421.4 \\
        Порог 100 км & 15.8 & 0.591 & 0.003 & 279.1 & 0.060 & 334.5 \\
        Экономическая & 10.0 & 0.011 & 0.066 & 278.7 & 0.058 & 282.4 \\
        \bottomrule
    \end{tabular}
    \tablesource
\end{table}

Пороговая матрица непригодна - 30 населенных пунктов не имеют соседей в радиусе 100 км. Для них ожидаемое
значение $y$ в SAR равно константе, а для остальных населенных пунктов - константе, деленной на
$(1 - \rho)$, поэтому при уровне $y \approx 18$ любое положительное $\rho$ резко ухудшает
правдоподобие, и оценка вырождается к нулю. Экономическая матрица не выявляет автокорреляции
($I = 0.011$), то есть сходство населенных пунктов по численности населения не воспроизводит обнаруженную
зависимость. Это согласуется с ее географическим характером, но не доказывает конкретный механизм.

\figuresingle{weights_profile.png}{Профиль правдоподобия SDM по масштабу затухания матрицы весов}{fig:profile}

\subsection*{Глобальная пространственная автокорреляция}

Все глобальные статистики отвергают гипотезу о случайном размещении населенных пунктов по выручке
(таблица \ref{tab:global}). Индекс Морана равен 0.582, и населенные пункты с выручкой выше средней в
основном окружены такими же населенными пунктами. Перестановочный тест дает тот же вывод. Индекс Гири меньше
единицы, что также соответствует положительной автокорреляции, а значимая статистика
Гетиса-Орда указывает на концентрацию высоких значений выручки. На диаграмме Морана 80\% населенных пунктов
лежат в квадрантах «высокий - высокий» (HH, 33.5\%) и «низкий - низкий» (LL, 46.5\%), а выбросы
«высокий - низкий» (HL) и «низкий - высокий» (LH) составляют 10.8\% и 9.1\% (рисунок
\ref{fig:moran_scatter}). Гипотеза 1 подтверждается.

\figuresingle{moran_scatter.png}{Диаграмма Морана для логарифма выручки, цветом выделены квадранты и подписаны их доли}{fig:moran_scatter}

\begin{table}[H]
    \centering
    \footnotesize
    \caption{Глобальные статистики пространственной автокорреляции}
    \label{tab:global}
    \begin{tabular}{lrrrr}
        \toprule
        \textbf{Статистика} & \textbf{Значение} & \textbf{Ожидание при $H_0$} & \textbf{$z$} & \textbf{$p$-value} \\
        \midrule
        Индекс Морана $I$ & 0.582 & -0.001 & 43.02 & $< 0.001$ \\
        Индекс Морана, 999 перестановок & 0.582 &  &  & 0.001 \\
        Индекс Гири $C$ & 0.434 & 1.000 & 38.76 & $< 0.001$ \\
        Статистика Гетиса-Орда $G$, 100 км & 0.02637 & 0.01814 & 18.64 & $< 0.001$ \\
        Двумерный $I_{xy}$ выручки и численности соседей & -0.020 &  &  & 0.038 \\
        Двумерный $I_{xy}$ выручки и автотрафика соседей & -0.021 &  &  & 0.042 \\
        \bottomrule
    \end{tabular}
    \par\vspace{1mm}{\footnotesize Примечание: для двумерных индексов приведено двустороннее псевдо
    $p$-value $2 (\min(m^-, m^+) + 1) / (R + 1)$ перестановочного теста $H_0: I_{xy} = 0$ против
    $H_1: I_{xy} \ne 0$, $R = 999$. Источник: расчеты авторов по данным X5 Group, Росстата и GeoNames.}
\end{table}

Двумерные индексы Морана отрицательны, но малы по величине. Населенные пункты, соседи которых крупнее и
загруженнее, в среднем имеют немного меньшую выручку на магазин. Среди регрессоров сильнее всего
пространственно сглажены показатели субъекта, индекс Морана равен 0.82 для реальной зарплаты и
0.65 для безработицы. Для численности населения он равен лишь 0.03, поэтому размер населенного пункта
пространственной структуры выручки не объясняет. Индекс Морана, рассчитанный по каждому из 26
месяцев, лежит в интервале от 0.51 до 0.61 (рисунок \ref{fig:monthly}).

\subsection*{Локальная пространственная автокорреляция}

Локальный индекс Морана с поправкой Бенджамини-Хохберга выделяет значимые кластеры у 484 населенных пунктов
из 933 (рисунок \ref{fig:local_moran}). Поправка Бенджамини-Иекутиели, верная при любой
зависимости статистик, оставляет 372 населенных пункта. На аналитических $p$-значениях FDR выделяет 467
населенных пунктов, а поправка Бонферрони - 254, поэтому различие между процедурами обусловлено прежде всего
строгостью поправки. Кластеры высокой выручки (224 населенных пункта) образуют Московская агломерация и
прилегающие области - Московская (84 населенных пункта), Ленинградская (30), Калужская (13), Тульская (12),
Белгородская и Воронежская (по 11), а также Мурманская область и Ямало-Ненецкий автономный округ.
Кластеры низкой выручки (212 населенных пунктов) сосредоточены в Краснодарском крае (38), Башкортостане (37),
Ставропольском крае (22), Ростовской (16) и Саратовской (15) областях. Выбросы «высокий - низкий»
(22 населенных пункта) чаще всего встречаются в Краснодарском крае (6 населенных пунктов). Границы кластеров во многом повторяют границы
субъектов, что указывает на роль региональных факторов.

\figuresingle{map_local_moran.png}{Кластеры локального индекса Морана, двусторонний перестановочный тест, FDR 5\%}{fig:local_moran}

Локальная статистика $G^*_i$ выделяет 246 населенных пунктов в горячих зонах и 241 в холодных, их
расположение совпадает с кластерами локального индекса Морана (рисунок \ref{fig:local_getis}).
Локальный индекс Гири выделяет 248 населенных пунктов, похожих на соседей при высоком уровне выручки, и 294
при низком (рисунок \ref{fig:local_geary}). По многомерному индексу Гири для выручки,
автотрафика и плотности сети 337 населенных пунктов входят в значимую область многомерного сходства с
соседями (рисунок \ref{fig:multivariate_geary}). Индекс агрегирует различия по трем признакам,
поэтому этот результат не означает значимого сходства по каждому показателю в отдельности. Локальный двумерный индекс Морана выручки и численности
соседей после поправки на множественное тестирование значимых населенных пунктов не выявляет.

\figurepair{map_local_getis.png}{Горячие и холодные зоны локальной статистики $G^*_i$, FDR 5\%}{fig:local_getis}
{map_local_geary.png}{Кластеры локального индекса Гири, FDR 5\%}{fig:local_geary}

\figuresingle{map_multivariate_geary.png}{Многомерный локальный индекс Гири по выручке, автотрафику и плотности сети, FDR 5\%}{fig:multivariate_geary}

\subsection*{Диагностика и выбор модели}

Оценки всех моделей приведены в таблице \ref{tab:models}, тесты ограничений - в таблице
\ref{tab:lr}. МНК-модель объясняет 63\% вариации выручки, однако ее остатки сохраняют сильную
автокорреляцию (индекс Морана 0.263, $z = 20.2$). Тесты множителей Лагранжа отвергают отсутствие
как лага ($LM_{lag} = 432.6$), так и пространственной ошибки ($LM_{err} = 373.6$), и значимы обе
робастные версии ($LM^{rob}_{lag} = 131.8$, $LM^{rob}_{err} = 72.7$). Большее значение робастного
теста на лаг указывает на SAR, но значимость обоих тестов говорит о том, что более общая модель
может быть предпочтительнее, поэтому выбор продолжается подходом от общей модели к частной.

В подходе от общей модели к частной ограничение $\lambda = 0$ в GNS не отвергается, и GNS
сводится к SDM, а все дальнейшие упрощения SDM (к SAR, к SLX и к SEM через ограничение общего
множителя) отвергаются. Расширение набора лагов $WX$ на все регрессоры модель не улучшает. Подход
от простой модели к сложной останавливается на SAR, подход от общей модели к частной приводит к
SDM, и LR-тест отвергает SAR в пользу SDM. У SDM наименьший AIC, и LM-тест не находит в ее
остатках автокорреляции, тогда как у SAR она значима на уровне 10\%. BIC сильнее штрафует шесть
дополнительных параметров и формально предпочитает SAR, однако LR-тест показывает, что эти
параметры значимы. В SARAR коэффициент $\lambda$ значим по тесту Вальда ($p = 0.03$), но LR-тест
ограничения $\lambda = 0$ его не отвергает ($p = 0.063$). Асимптотически эквивалентные тесты
расходятся в конечной выборке, что указывает на слабую пространственную зависимость в ошибке после
учета лага. Для дальнейшей интерпретации выбрана SDM.

\begin{table}[H]
    \centering
    \footnotesize
    \caption{Сравнение пространственных моделей}
    \label{tab:models}
    \begin{tabular}{lrrrrrrl}
        \toprule
        \textbf{Модель} & \textbf{$\rho$} & \textbf{$\lambda$} & \textbf{Парам.} & \textbf{$\ln L$} & \textbf{AIC} & \textbf{BIC} & \textbf{Остат. автокорр.} \\
        \midrule
        МНК &  &  & 15 & 277.9 & -525.8 & -453.2 & 0.263*** \\
        SLX &  &  & 21 & 349.8 & -657.6 & -556.0 & 0.162*** \\
        SAR & 0.540*** &  & 16 & 417.1 & -802.1 & -724.7 & LM = 2.84* \\
        SEM &  & 0.747*** & 16 & 407.0 & -782.0 & -704.6 &  \\
        SDM & 0.474*** &  & 22 & 432.3 & -820.7 & -714.2 & LM = 1.15 \\
        SDEM &  & 0.686*** & 22 & 414.4 & -784.8 & -678.4 &  \\
        SARAR & 0.488*** & 0.195** & 17 & 418.8 & -803.6 & -721.3 &  \\
        GNS & 0.446*** & 0.110 & 23 & 432.9 & -819.9 & -708.6 &  \\
        \bottomrule
    \end{tabular}
    \par\vspace{1mm}{\footnotesize Примечание: *, **, *** - значимость на уровнях 10\%, 5\% и 1\%. В последнем столбце для МНК и SLX приведен индекс Морана остатков, для SAR и SDM - LM-тест остаточной автокорреляции. Источник: расчеты авторов по данным X5 Group, Росстата и GeoNames.}
\end{table}

\begin{table}[H]
    \centering
    \footnotesize
    \caption{Тесты отношения правдоподобия для вложенных моделей}
    \label{tab:lr}
    \begin{tabular}{llrrr}
        \toprule
        \textbf{Переход} & \textbf{Ограничение} & \textbf{LR} & \textbf{$q$} & \textbf{$p$-value} \\
        \midrule
        GNS $\to$ SDM & $\lambda = 0$ & 1.21 & 1 & 0.271 \\
        GNS $\to$ SDEM & $\rho = 0$ & 37.08 & 1 & $< 0.001$ \\
        GNS $\to$ SARAR & $\theta = 0$ & 28.30 & 6 & $< 0.001$ \\
        SDM $\to$ SAR & $\theta = 0$ & 30.55 & 6 & $< 0.001$ \\
        SDM, все $WX$ $\to$ SEM & $\theta = -\rho \beta$ & 60.52 & 13 & $< 0.001$ \\
        SDM, все $WX$ $\to$ SDM & 7 лагов равны нулю & 9.87 & 7 & 0.196 \\
        SDM $\to$ SLX & $\rho = 0$ & 165.04 & 1 & $< 0.001$ \\
        SDEM $\to$ SEM & $\theta = 0$ & 14.78 & 6 & 0.022 \\
        SARAR $\to$ SAR & $\lambda = 0$ & 3.45 & 1 & 0.063 \\
        SAR $\to$ МНК & $\rho = 0$ & 278.35 & 1 & $< 0.001$ \\
        SDM $\to$ МНК & $\rho = 0$, $\theta = 0$ & 308.90 & 7 & $< 0.001$ \\
        \bottomrule
    \end{tabular}
    \tablesource
\end{table}

Тест Бройша-Пагана выявляет гетероскедастичность в SDM ($BP = 70.1$, $p < 0.001$), поэтому
стандартные ошибки максимального правдоподобия и LR-тесты, построенные при однородной дисперсии,
не вполне надежны. Выбор SDM по AIC и правдоподобию отражает качество подгонки, а формальный
вывод о ключевом ограничении проверяется робастно. Тест Вальда с ковариационной матрицей
двухшагового МНК, устойчивой к гетероскедастичности, отвергает $\theta = 0$ ($W = 23.7$, $q = 6$,
$p < 0.001$). Пространственный двухшаговый МНК
со стандартными ошибками Уайта дает $\hat{\rho} = 0.495$ (стандартная ошибка 0.052) против 0.474
у метода максимального правдоподобия (таблица \ref{tab:sdm}). Знаки всех коэффициентов и
значимость всех, кроме лага автомобильного трафика ($p = 0.11$), сохраняются. Двухшаговый МНК для
SAR дает $\hat{\rho} = 0.567$ против 0.540 у метода максимального правдоподобия. Обобщенный МНК
Келеджиана-Пручи для SEM дает $\hat{\lambda} = 0.623$ со стандартной ошибкой 0.146, что заметно
ниже оценки максимального правдоподобия (0.747), но эта оценка не робастна к гетероскедастичности.
Робастный GMM для SEM дает $\hat{\lambda} = 0.717$ (стандартная ошибка 0.032), что близко к оценке
максимального правдоподобия.

\begin{table}[H]
    \centering
    \footnotesize
    \caption{Оценки коэффициентов МНК и SDM}
    \label{tab:sdm}
    \begin{tabular}{lrrr}
        \toprule
        \textbf{Регрессор} & \textbf{МНК} & \textbf{SDM, ML} & \textbf{SDM, 2МНК робастный} \\
        \midrule
        Численность населения & -0.022*** (0.006) & -0.001 (0.006) & 0.000 (0.006) \\
        Домохозяйства & 0.046*** (0.014) & 0.037*** (0.011) & 0.037*** (0.012) \\
        Пеший трафик & 0.006 (0.019) & 0.015 (0.018) & 0.016 (0.019) \\
        Автомобильный трафик & 0.048*** (0.012) & 0.024** (0.009) & 0.023** (0.010) \\
        Продуктовые конкуренты & -0.035*** (0.006) & -0.028*** (0.005) & -0.028*** (0.006) \\
        Плотность сети & 0.035** (0.015) & -0.003 (0.014) & -0.003 (0.014) \\
        Пункты выдачи маркетплейсов & 0.031*** (0.008) & 0.027*** (0.007) & 0.027*** (0.007) \\
        Медицинские учреждения и аптеки & 0.027*** (0.007) & 0.011* (0.006) & 0.011* (0.006) \\
        Кассы & 0.099*** (0.006) & 0.078*** (0.004) & 0.077*** (0.006) \\
        Доля больших магазинов & 0.218*** (0.040) & 0.155*** (0.031) & 0.153*** (0.034) \\
        Доля новых магазинов & -0.026 (0.031) & 0.060** (0.026) & 0.061** (0.028) \\
        Реальная зарплата & 0.413*** (0.033) & 0.207*** (0.039) & 0.193*** (0.042) \\
        Безработица & 0.034*** (0.010) & 0.026*** (0.007) & 0.026*** (0.008) \\
        $W \cdot$ численность населения &  & -0.059*** (0.019) & -0.055** (0.024) \\
        $W \cdot$ домохозяйства &  & -0.018 (0.030) & -0.023 (0.031) \\
        $W \cdot$ пеший трафик &  & 0.002 (0.035) & 0.000 (0.038) \\
        $W \cdot$ автомобильный трафик &  & 0.054** (0.026) & 0.047 (0.030) \\
        $W \cdot$ продуктовые конкуренты &  & 0.026** (0.012) & 0.027** (0.012) \\
        $W \cdot$ плотность сети &  & 0.084** (0.035) & 0.080** (0.035) \\
        $\rho$ &  & 0.474*** (0.033) & 0.495*** (0.052) \\
        \midrule
        Наблюдений & 933 & 933 & 933 \\
        $R^2$ / псевдо-$R^2$ Нагелькерке & 0.629 & 0.733 & \\
        \bottomrule
    \end{tabular}
    \par\vspace{1mm}{\footnotesize Примечание: в скобках стандартные ошибки; *, **, *** - значимость на уровнях 10\%, 5\% и 1\%; для МНК использованы стандартные ошибки, робастные к гетероскедастичности (HC1), для SDM ML - асимптотические стандартные ошибки максимального правдоподобия, для 2МНК - стандартные ошибки Уайта. Источник: расчеты авторов по данным X5 Group, Росстата и GeoNames.}
\end{table}

Знаки некоторых коэффициентов не соответствуют ожидаемому направлению эффектов регрессоров:

\begin{enumerate}
    \item Пункты выдачи маркетплейсов связаны с большей, а не меньшей выручкой. Возможны два
    объяснения. Во-первых, пункты выдачи могут отмечать плотную городскую застройку с высоким
    пешеходным потоком. Во-вторых, на маркетплейсах в основном покупают непродовольственные товары, поэтому
    пункты выдачи слабо отбирают спрос у продуктового магазина, а покупатель, пришедший за заказом,
    может заодно зайти в магазин. Данные не позволяют проверить и разделить эти механизмы;
    \item Положительный коэффициент безработицы отражает межрегиональные различия, не учтенные
    зарплатой, и в спецификации с эффектами субъектов этот показатель поглощается;
    \item Доля новых магазинов меняет знак с отрицательного в МНК на положительный в SDM. Возможное
    объяснение, которое в работе не проверялось, состоит в том, что новые магазины чаще
    открываются в регионах с низкой выручкой, и после учета выручки соседей их собственная выручка
    оказывается выше;
    \item Отрицательный прямой эффект численности населения в МНК в SDM исчезает и переходит в
    косвенный. Возможное объяснение состоит в том, что в крупных населенных пунктах и рядом с ними
    сеть плотнее и выручка на один магазин ниже. Это объяснение остается гипотезой, поскольку
    число магазинов сети в населенном пункте в основную модель не входит, а его добавление в
    проверке устойчивости эту связь не меняет.
\end{enumerate}

\subsection*{Прямые, косвенные и общие эффекты}

Коэффициент пространственного лага $\hat{\rho} = 0.474$ положителен и значим, гипотеза 2
подтверждается, и после учета характеристик населенных пунктов выручка соседних рынков изменяется
однонаправленно. Сам коэффициент $\rho$ не является предельным эффектом; он определяет
мультипликатор $1 / (1 - \hat{\rho}) \approx 1.9$, во сколько раз общий эффект регрессора без
собственного лага превышает его коэффициент при одновременном изменении во всех населенных пунктах. Слагаемые
ряда $(\rho W)^q (\beta I + \theta W)$ убывают примерно вдвое с каждым $q$ (рисунок
\ref{fig:orders}); слагаемое при $q = 0$ уже включает $\theta W$, то есть влияние прямых соседей
через лаги регрессоров.

Поскольку зависимая переменная - логарифм выручки, эффекты логарифмированных регрессоров
интерпретируются как эластичности, а эффекты остальных - как полуэластичности. Эффекты SDM
приведены в таблице \ref{tab:impacts}; в двух последних столбцах указаны $p$-значения прямых и
косвенных эффектов по робастному двухшаговому МНК.

\begin{table}[H]
    \centering
    \footnotesize
    \caption{Прямые, косвенные и общие эффекты SDM}
    \label{tab:impacts}
    \begin{tabular}{lrrrrr}
        \toprule
        \multirow{2}{*}{\textbf{Регрессор}} & \multicolumn{3}{c}{\textbf{ML}} & \multicolumn{2}{c}{\textbf{$p$, 2МНК робастный}} \\
        & \textbf{Прямой} & \textbf{Косвенный} & \textbf{Общий} & \textbf{Прямой} & \textbf{Косвенный} \\
        \midrule
        Численность населения & -0.004 & -0.110*** & -0.115*** & 0.594 & 0.013 \\
        Домохозяйства & 0.037*** & -0.001 & 0.037 & 0.002 & 0.828 \\
        Пеший трафик & 0.016 & 0.017 & 0.032 & 0.397 & 0.780 \\
        Автомобильный трафик & 0.028*** & 0.121** & 0.149*** & 0.006 & 0.044 \\
        Продуктовые конкуренты & -0.027*** & 0.024 & -0.003 & $< 0.001$ & 0.233 \\
        Плотность сети & 0.002 & 0.153** & 0.155** & 0.949 & 0.013 \\
        Пункты выдачи маркетплейсов & 0.028*** & 0.024*** & 0.052*** & $< 0.001$ & $< 0.001$ \\
        Медицинские учреждения и аптеки & 0.012** & 0.010** & 0.021** & 0.049 & 0.065 \\
        Кассы & 0.080*** & 0.068*** & 0.148*** & $< 0.001$ & $< 0.001$ \\
        Доля больших магазинов & 0.159*** & 0.136*** & 0.294*** & $< 0.001$ & $< 0.001$ \\
        Доля новых магазинов & 0.062** & 0.053** & 0.114** & 0.030 & 0.059 \\
        Реальная зарплата & 0.212*** & 0.181*** & 0.393*** & $< 0.001$ & $< 0.001$ \\
        Безработица & 0.027*** & 0.023*** & 0.050*** & 0.002 & 0.002 \\
        \bottomrule
    \end{tabular}
    \par\vspace{1mm}{\footnotesize Примечание: значимость по 2000 симуляциям; *, **, *** - на уровнях 10\%, 5\% и 1\%. Источник: расчеты авторов по данным X5 Group, Росстата и GeoNames.}
\end{table}

Рост среднего по магазинам населенного пункта числа продуктовых конкурентов в радиусе 500 м на единицу связан
с выручкой магазина ниже на 2.66\% ($\exp(-0.027) - 1$, $p < 0.001$ в обеих оценках), и гипотеза
3 подтверждается в части конкурентов. Регрессор является средним счетчиком по магазинам населенного пункта,
поэтому оценка описывает среднюю условную связь на уровне населенного пункта, а не последствия
открытия одного конкурента рядом с конкретным магазином. Новый магазин может попасть в радиусы
нескольких магазинов сети и по-разному изменить средний счетчик. Результат также описывает связь,
а не причинное влияние, поскольку конкуренты тоже выбирают места размещения.
По 2000 симуляциям 95\%-интервал прямого эффекта конкурентов составляет от $-0.038$ до $-0.017$.
Косвенный эффект конкурентов незначим, поэтому общий эффект близок к нулю и содержательно не
интерпретируется. Прямой эффект плотности собственной сети статистически не отличается от нуля
(0.002, $p = 0.89$, 95\%-интервал от $-0.025$ до $0.028$). Нулевая гипотеза не отвергается, то есть
свидетельств каннибализации недостаточно, однако интервал не исключает эффектов порядка
$\pm 2.5\%$. Этот вывод подвержен той же эндогенности размещения.

В основной спецификации значимы косвенные эффекты трех показателей спроса и конкуренции. Косвенный
эффект численности населения отрицателен ($-0.110$), и населенные пункты, соседи которых крупнее, имеют меньшую
выручку на магазин, что согласуется с оттоком покупателей в крупные центры. Косвенный эффект
автомобильного трафика положителен ($0.121$), что согласуется с агломерационным механизмом
транспортных коридоров. Косвенный эффект плотности сети положителен ($0.153$), что согласуется с
выгодами плотности сети \cite{holmes2011}, но может отражать и ненаблюдаемый спрос в окрестности.
Для регрессоров без собственных лагов $WX$ (кассы, зарплата, безработица и другие) косвенный
эффект пропорционален прямому с множителем, заданным $\rho$, и как самостоятельный спилловер не
интерпретируется.

\figurepair{impacts_by_order.png}{Слагаемые ряда мультипликатора в общем эффекте, SDM}{fig:orders}
{impacts_comparison.png}{Косвенные эффекты ключевых переменных в разных моделях с 95\%-интервалами}{fig:indirect}

Локальные спилловеры SLX и SDEM, затрагивающие только прямых соседей, имеют те же знаки для
численности, автотрафика и плотности сети, а для конкурентов знак различается (0.024 в SDM и
$-0.006$ в SDEM, оба незначимы) (рисунок \ref{fig:indirect}). В SLX со стандартными ошибками, робастными к гетероскедастичности,
значимы лаги численности ($-0.158$), автотрафика ($0.206$) и плотности сети ($0.164$), однако в
остатках SLX сохраняется автокорреляция (индекс Морана 0.162), и их значимость завышена. В SDEM
значимы косвенные эффекты численности ($-0.067$, $p = 0.008$) и плотности сети ($0.154$,
$p = 0.004$). В SAR косвенные эффекты по построению пропорциональны прямым, поэтому эта модель
приписывает соседям те же знаки, что и самому населенному пункту, и не может разделить внутренний и
межрыночный механизмы.

Гипотеза 5 подтверждается, совместный тест $\rho = 0$ и $\theta = 0$ отвергается ($LR = 308.9$,
$q = 7$, $p < 0.001$), AIC снижается с -525.8 у МНК до -820.7 у SDM, а остаточная автокорреляция
исчезает.

\subsection*{Региональная неоднородность и устойчивость результатов}

Главная проверка касается региональных различий, на которые указывают доля дисперсии, объясняемая
субъектом, и совпадение кластеров LISA с границами субъектов. Кластеризация стандартных ошибок МНК
по 67 субъектам, учитывающая общий источник вариации зарплаты и безработицы, сохраняет значимость
зарплаты ($p < 0.001$), а значимость безработицы снижается до $p = 0.026$. В SDM с фиксированными эффектами субъектов, которые
поглощают региональные уровни цен, доходов и особенности управления сетью, коэффициент $\rho$
снижается с 0.474 до 0.204, но остается значимым (таблица \ref{tab:robustness}). Индекс Морана
остатков МНК с эффектами субъектов равен 0.031 при ожидаемом -0.042 ($p < 0.001$). Внутри субъектов
пространственная связь выручки сохраняется, но слабее. Косвенные эффекты численности,
автотрафика и плотности сети в этой спецификации незначимы. Следовательно, косвенные эффекты
основной модели чувствительны к учету региональных факторов и не могут считаться установленным
взаимодействием соседних рынков. Потеря значимости может отражать устранение смешения с
региональными факторами, уменьшение полезной вариации или рост неопределенности, а различие
«значим - незначим» само по себе не является тестом различия эффектов. Отдельно межрегиональная
компонента не оценивалась, поэтому гипотеза 4 подтверждается только в основной спецификации, и
этот результат неустойчив к эффектам субъектов.

\begin{table}[H]
    \centering
    \footnotesize
    \setlength{\tabcolsep}{4pt}
    \caption{Проверка устойчивости SDM - $\rho$, прямой эффект конкурентов и косвенные эффекты}
    \label{tab:robustness}
    \begin{tabular}{>{\raggedright\arraybackslash}p{4.6cm}rrrrrrr}
        \toprule
        & & & \textbf{ADE} & \multicolumn{4}{c}{\textbf{AIE}} \\
        \cmidrule(lr){4-4} \cmidrule(lr){5-8}
        \textbf{Вариант} & \textbf{$N$} & \textbf{$\rho$} & \textbf{Конкур.} & \textbf{Числен.} & \textbf{Автотраф.} & \textbf{Конкур.} & \textbf{Плотн.} \\
        \midrule
        Основная спецификация & 933 & 0.474*** & -0.027*** & -0.110*** & 0.121** & 0.024 & 0.153** \\
        Без касс и доли больших магазинов & 933 & 0.619*** & -0.033*** & -0.167*** & 0.306*** & -0.017 & 0.351*** \\
        Эффекты субъектов & 933 & 0.204*** & -0.026*** & -0.004 & 0.008 & 0.004 & 0.075 \\
        Эффекты субъектов, без касс и доли больших & 933 & 0.280*** & -0.029*** & 0.000 & 0.110* & 0.005 & 0.147 \\
        \midrule
        kNN, $k = 8$ & 933 & 0.467*** & -0.026*** & -0.114*** & 0.151*** & 0.021 & 0.107** \\
        $1/d$, 200 км & 933 & 0.475*** & -0.027*** & -0.096*** & 0.109** & 0.024 & 0.142** \\
        $1/d^2$, 400 км & 933 & 0.461*** & -0.024*** & -0.029 & 0.075** & -0.005 & 0.179*** \\
        Экономическая матрица & 933 & 0.058 & -0.033*** & 0.076 & -0.006 & -0.018 & 0.071 \\
        \midrule
        5 и более магазинов & 570 & 0.533*** & -0.022*** & -0.048 & -0.051 & 0.042* & 0.096 \\
        Без Москвы и Санкт-Петербурга & 931 & 0.471*** & -0.027*** & -0.111*** & 0.121*** & 0.023 & 0.152** \\
        Численность по данным Росстата & 822 & 0.457*** & -0.024*** & -0.125*** & 0.098** & 0.046** & 0.104* \\
        Без магазинов со скачками выручки & 841 & 0.452*** & -0.029*** & -0.153*** & 0.117*** & 0.023 & 0.087 \\
        Без населенных пунктов без соседей в 100 км & 903 & 0.494*** & -0.022*** & -0.089** & 0.150*** & 0.012 & 0.194*** \\
        \midrule
        Медианная выручка магазина & 933 & 0.445*** & -0.030*** & -0.112*** & 0.115** & 0.034 & 0.124* \\
        Окно 24 месяца & 933 & 0.475*** & -0.027*** & -0.113*** & 0.126*** & 0.024 & 0.156** \\
        С числом магазинов в населенном пункте & 933 & 0.475*** & -0.027*** & -0.110*** & 0.123*** & 0.024 & 0.156** \\
        \bottomrule
    \end{tabular}
    \par\vspace{1mm}{\footnotesize Примечание: ADE - прямой эффект, AIE - косвенный эффект; *, **, *** - значимость на уровнях 10\%, 5\% и 1\% по 2000 симуляциям. Источник: расчеты авторов по данным X5 Group, Росстата и GeoNames.}
\end{table}

Исключение касс и доли больших магазинов повышает $\rho$ до 0.619 и увеличивает косвенные эффекты
автотрафика и плотности сети в 2-2.5 раза. Это показывает чувствительность оценок к набору
контролей формата, но не доказывает, что основная спецификация дает нижнюю границу истинного
эффекта, поскольку контроли эндогенны и направление смещения заранее неизвестно. На географических матрицах $\rho$ находится в пределах 0.46-0.48, но эффект
численности незначим на матрице обратных квадратов расстояний, а эффект плотности сети - в
нескольких подвыборках. На строгой подвыборке из 570 населенных пунктов с пятью и более магазинами
косвенный эффект автотрафика меняет знак, а эффекты численности и плотности сети незначимы. Эта
подвыборка охватывает все 67 субъектов, но состоит из более крупных населенных пунктов (медиана численности
45 тыс. человек против 27 тыс. в полной выборке), поэтому она одновременно меняет состав рынков и
матрицу весов и не изолирует эффект усреднения шума. Исключение 30 населенных пунктов без соседей в радиусе
100 км, окно ровно из 24 месяцев и добавление логарифма числа магазинов сети в населенном пункте основные
оценки почти не меняют; коэффициент при числе магазинов незначим ($-0.005$, $p = 0.74$), что не
подтверждает объяснение отрицательной связи с численностью соседей через распределение спроса
между магазинами сети, но и не исключает его, поскольку число магазинов выбирает сеть. С
учетом шести проверяемых косвенных эффектов поправку Бонферрони ($p < 0.0083$) в основной
спецификации проходит только эффект численности, а по робастным $p$-значениям - ни один. На
экономической матрице по численности населения эффекты пропадают, то есть эта альтернативная
структура связей их не воспроизводит. Прямой эффект конкурентов во всех 16 вариантах значим на
уровне 1\% и лежит в пределах от $-0.033$ до $-0.022$, то есть от $-3.3\%$ до $-2.2\%$.

Сила пространственной зависимости во времени сопоставима. Оцененный по каждому месяцу коэффициент
$\rho$ SDM лежит в пределах 0.45-0.53, индекс Морана - в пределах 0.51-0.61 (рисунок
\ref{fig:monthly}).

\subsection*{Пространственная панель}

Тест Хаусмана отвергает модель со случайными эффектами ($\chi^2 = 267$, $p < 0.001$), поэтому в
панели используются двунаправленные фиксированные эффекты. После их учета остатки сохраняют
зависимость между населенными пунктами, CD-тест Песарана равен 47.2 ($p < 0.001$), средняя корреляция
остатков соседей - 0.111. Панельные LM-тесты отвергают отсутствие как лага ($LM = 8137$), так и
пространственной ошибки ($LM = 8162$), а робастные версии указывают на ошибку (48.2 против 23.7).
Байесовские апостериорные вероятности дают менее однозначную картину. В статической постановке
93.5\% приходится на SDM, в динамической 76.4\% на SDEM и 23.6\% на SDM. Общим для всех критериев
является сильная пространственная зависимость, оценки $\rho$ и $\lambda$ близки и составляют 0.67
(таблица \ref{tab:panel}).

\begin{table}[H]
    \centering
    \footnotesize
    \caption{Пространственные панельные модели с двунаправленными фиксированными эффектами, $N = 933$, $T = 26$}
    \label{tab:panel}
    \begin{tabular}{lrrrrr}
        \toprule
        & \textbf{FE} & \textbf{SAR FE} & \textbf{SEM FE} & \textbf{Динам. SAR} & \textbf{Динам. SDM} \\
        \midrule
        $y_{t-1}$ & & & & 0.831*** & 0.831*** \\
        & & & & (0.004) & (0.004) \\
        $W y_{t}$ ($\rho$) & & 0.671*** & & 0.643*** & 0.651*** \\
        & & (0.008) & & (0.009) & (0.009) \\
        $W y_{t-1}$ ($\eta$) & & & & -0.520*** & -0.526*** \\
        & & & & (0.012) & (0.012) \\
        $W u_t$ ($\lambda$) & & & 0.672*** & & \\
        & & & (0.008) & & \\
        Реальная зарплата & 0.053 & 0.038** & 0.139*** & -0.019* & -0.141*** \\
        & (0.083) & (0.015) & (0.027) & (0.010) & (0.022) \\
        Безработица & 0.003 & 0.001 & 0.001 & -0.002** & -0.003** \\
        & (0.006) & (0.001) & (0.002) & (0.001) & (0.002) \\
        \bottomrule
    \end{tabular}
    \par\vspace{1mm}{\footnotesize Примечание: в скобках стандартные ошибки; *, **, *** - значимость на уровнях 10\%, 5\% и 1\%. Для FE использованы стандартные ошибки Дрисколла-Краая, устойчивые к сериальной и пространственной корреляции. Для остальных моделей приведены стандартные ошибки максимального правдоподобия, которые предполагают независимость ошибок во времени, поэтому их значимость условна относительно этой предпосылки. Динамические модели оценены с коррекцией смещения Ли-Ю; в динамической SDM лаг зарплаты равен 0.136 (0.027), лаг безработицы незначим. Источник: расчеты авторов по данным X5 Group, Росстата и GeoNames.}
\end{table}

В остатках пространственных моделей SAR FE и SEM FE средний по месяцам индекс Морана близок к
нулю ($-0.002$ и $-0.001$ против 0.242 у модели FE без пространственных членов). Модели с лагами
региональных регрессоров отвергают ограничения SAR и SEM. LR-тест SDM FE против SAR FE равен 43.5,
SDEM FE против SEM FE - 14.0 (оба $p < 0.001$). В модели SARAR FE коэффициент при пространственной
ошибке (0.447, $p = 0.009$) значимее коэффициента при лаге (0.357, $p = 0.058$), что согласуется с
робастными LM-тестами. Тест Вулдриджа отвергает отсутствие сериальной корреляции ошибок модели FE
($\chi^2 = 15651$, $p < 0.001$). Со стандартными ошибками Дрисколла-Краая \cite{driscoll1998} коэффициент при реальной
зарплате (0.053) незначим ($p = 0.52$), и в таблице \ref{tab:panel} для FE приведены именно они.
Для пространственных моделей эта поправка не определена, поэтому проверена остаточная динамика.
Корреляция остатка населенного пункта с его остатком в предыдущем месяце равна 0.76 у FE, $-0.02$ у SAR FE и
0.15 у SEM FE. Пространственный лаг поглощает основную часть сериальной корреляции, но в SEM FE она
сохраняется, а стандартные ошибки максимального правдоподобия ее не учитывают. Поэтому звезды
SAR FE и SEM FE и основанные на них LR-тесты и сравнения моделей условны относительно предпосылок
метода максимального правдоподобия, а динамика учитывается явно в динамических моделях. Во времени в панели меняются только два региональных
показателя, поэтому гипотезы 3 и 4 о характеристиках магазинов и соседних рынков в панели не
проверяются.

Сильная зависимость месячных отклонений может отражать шоки двух видов - общие для всего субъекта
(региональная сезонность, курортный спрос, изменения доходов) и собственно локальные. Эффекты
«субъект $\times$ месяц» поглощают первые. После их учета средний индекс Морана остатков по месяцам
снижается с 0.17-0.35 до значений от $-0.03$ до 0.08. Обычный тест Морана не учитывает, что
остатки получены после оценки эффектов, поэтому его $p$-значения приближенные. На 5\%-м уровне
индекс значим в 14 месяцах из 26, после поправки FDR на 26 проверок - в 12. CD-тест Песарана снижается с 47.2 до 26.3 и по-прежнему отвергает отсутствие
зависимости между населенными пунктами. Локальная автокорреляция остатков после учета региональных шоков
мала, но общая зависимость между населенными пунктами сохраняется.

Гипотеза 6 подтверждается для обеих моделей при предпосылках метода максимального правдоподобия.
При двунаправленных фиксированных эффектах $\rho$ в SAR FE и $\lambda$ в SEM FE значимы, а их
$z$-статистики около 89, и вывод изменился бы, только если бы стандартные ошибки были занижены
более чем в 45 раз. Учет региональных месячных шоков показывает, что большая
часть этой связи приходится на общую динамику населенных пунктов одного субъекта.

\subsection*{Динамическая панель}

Устойчивость месячной выручки оценивается несколькими методами (таблица \ref{tab:dynamic}). Для
сопоставимости МНК и модель с фиксированными эффектами оценены на тех же месяцах, что и GMM с
лагами 2-4 (22 392 наблюдения), и с теми же эффектами времени. МНК дает 0.983, оценка с
фиксированными эффектами - 0.758, и оценка Ареллано-Бонда 0.817 лежит между ними. Это
согласуется с ожидаемыми направлениями смещений, хотя положение GMM-оценки между МНК и FE не
является теоремой для любых спецификаций, особенно при зависимости ошибок между населенными пунктами.
Смещение Никелла при $T = 26$ имеет порядок $(1 + \gamma) / (T - 1) \approx 0.07$, около 9\% от
$\gamma$. Тест AR(2) не отвергает отсутствие автокорреляции второго порядка ($p = 0.23$), но это не
подтверждает остальные предпосылки метода. При инструментах с третьего лага оценка снижается до
0.768; эта оценка получена на меньшей выборке (21 459 наблюдений), поэтому сдвиг без отдельной
диагностики не доказывает слабость инструментов. Системный GMM дает 0.958, что также лежит внутри
интервала между FE и МНК.

$J$-тест Хансена-Саргана отвергает моментные условия во всех трех спецификациях ($p < 0.001$; в
модели Ареллано-Бонда 29 инструментов при 27 параметрах, включая эффекты времени, то есть два
ограничения). Возможных причин несколько. Во-первых, зависимость ошибок между населенными пунктами, на которую
указывает CD-тест, может нарушать калибровку теста, но сама по себе не тождественна нарушению
моментных условий. Во-вторых, возможны пропущенная динамика и нарушение строгой экзогенности зарплаты и
безработицы, которые в модели служат инструментами для самих себя. Поэтому GMM-оценки не
используются как подтвержденная оценка инерции. При $\hat{\gamma} = 0.817$ долгосрочный эффект
регрессора был бы в 5.5 раза больше краткосрочного, однако это условная иллюстрация, а
коэффициент при реальной зарплате в модели Ареллано-Бонда незначим (0.031, $p = 0.11$).
Региональные регрессоры почти не меняются от месяца к месяцу, поэтому их коэффициенты в
динамических моделях неустойчивы по знаку. В модели FE коэффициент при зарплате отрицателен
($-0.073$), а в модели Ареллано-Бонда коэффициент при безработице положителен (0.011).

\begin{table}[H]
    \centering
    \footnotesize
    \caption{Оценки коэффициента при $y_{t-1}$ в динамической панели, $N = 933$, $T = 26$}
    \label{tab:dynamic}
    \begin{tabular}{lrrrrr}
        \toprule
        \textbf{Оценка} & \textbf{$\hat{\gamma}$} & \textbf{Ст. ошибка} & \textbf{$J$, $p$ (df)} & \textbf{AR(1), $p$} & \textbf{AR(2), $p$} \\
        \midrule
        МНК с эффектами времени & 0.983 & 0.002 & & & \\
        Фиксированные эффекты & 0.758 & 0.008 & & & \\
        Ареллано-Бонд, лаги 2-4 & 0.817 & 0.034 & $< 0.001$ (2) & $< 0.001$ & 0.226 \\
        Системный GMM, лаги 2-4 & 0.958 & 0.021 & $< 0.001$ (5) & $< 0.001$ & 0.150 \\
        Ареллано-Бонд, лаги 3-5 & 0.768 & 0.037 & $< 0.001$ (2) & $< 0.001$ & 0.184 \\
        Динамическая SAR с FE & 0.831 & 0.004 & & & \\
        \bottomrule
    \end{tabular}
    \par\vspace{1mm}{\footnotesize Примечание: МНК и FE оценены на месяцах 3-26, как GMM с лагами 2-4; стандартные ошибки МНК и FE кластеризованы по населенным пунктам, GMM - робастные двухшаговые, у динамической SAR - стандартные ошибки максимального правдоподобия после коррекции смещения Ли-Ю; динамическая SAR дополнительно содержит $\rho W y_t$ и $\eta W y_{t-1}$, поэтому ее $\hat{\gamma}$ сопоставима с остальными оценками лишь приближенно. Источник: расчеты авторов по данным X5 Group, Росстата и GeoNames.}
\end{table}

Пространственная динамическая модель с коррекцией смещения Ли-Ю показывает высокую устойчивость
месячной выручки ($\hat{\gamma} = 0.831$ против 0.760 без коррекции) и сильную одновременную
связь с соседями ($\hat{\rho} = 0.643$). Коррекция повышает $\hat{\gamma}$ на 0.07, что совпадает
по порядку с ожидаемым смещением $1/T$. Лаг выручки соседей в прошлом месяце отрицателен
($\hat{\eta} = -0.520$) и близок к $-\hat{\gamma} \hat{\rho} = -0.534$. Тест Вальда для
ограничения $H_0: \eta + \gamma \rho = 0$ по скорректированным оценкам, рассчитанный
дельта-методом, дает разность 0.015 со стандартной ошибкой 0.008 ($W = 3.37$, $p = 0.066$). На
5\%-м уровне ограничение не отвергается, на 10\%-м отвергается, поэтому оно не подтверждается, а
лишь не отвергается при заранее выбранном уровне 5\%. При этом ограничении исчезает пространственная
передача через прошлые значения выручки, но, как показывает преобразованное уравнение в разделе
о модели, пространственный мультипликатор $(I - \rho W)^{-1}$ для текущих регрессоров и шоков
сохраняется. Данные поэтому согласуются с интерпретацией, в которой месячная динамика выручки
соседей связана прежде всего через одновременные пространственно коррелированные шоки, а не через
передачу прошлой выручки от населенного пункта к населенному пункту, хотя строго ее не доказывают. С этим согласуются
робастные LM-тесты, модель SARAR FE и байесовские вероятности, отдающие предпочтение SDEM в
динамике; эти критерии построены при предпосылках метода максимального правдоподобия. Динамическая
SDEM в использованной библиотеке не реализована, поэтому оценены динамические SAR и SDM. В
динамической SDM коэффициент при зарплате отрицателен ($-0.141$), а ее пространственный лаг
положителен (0.136) и близок к нему по модулю, поэтому эффекты месячной вариации региональной
зарплаты в динамике содержательно не интерпретируются.

\subsection*{Неоднородность эффектов, географически взвешенная регрессия}

Кросс-валидация выбирает адаптивное окно в 61 ближайший населенный пункт (6.5\% выборки), эффективное число
параметров модели $\operatorname{tr}(S)$ равно 109.6, квази-$R^2$ - 0.765. Модели сравниваются по
одному критерию AICc, учитывающему размер выборки. У МНК он равен $-525.3$, у GWR $-730.0$, у SDM
$-819.6$. GWR лучше МНК, но уступает SDM, а в ее остатках сохраняется пространственная
автокорреляция (индекс Морана 0.102, обычный тест не учитывает локальное оценивание, поэтому его
$p < 0.001$ приближенное), тогда как в остатках SDM ее нет. Преимущество SDM по AICc показывает
сравнительное качество этих спецификаций, но не доказывает, что зависимость является истинной
причиной паттерна, а неоднородность коэффициентов отсутствует.
Поскольку в GWR нет пространственного лага, ее локальные коэффициенты сопоставимы с оценками МНК,
а не SDM.

Остатки GWR пространственно зависимы, поэтому локальные $t$-тесты и F-тест предполагают неверную
структуру ошибок, поэтому приведенные ниже выводы о статистической значимости следует считать
предварительными.
Глобальный VIF не проверяет локальную коллинеарность, которая особенно вероятна для зарплаты и
безработицы, одинаковых внутри субъекта. F-тест Ю. Леунга, Ч.-Л. Ми и В.-С. Чжана \cite{leung2000} ($H_0$: коэффициент одинаков
во всех населенных пунктах) отвергает постоянство
коэффициентов при продуктовых конкурентах, автомобильном и пешем трафике, плотности сети, доле
новых магазинов и безработице. Для численности населения, домохозяйств, касс, доли больших
магазинов и реальной зарплаты неоднородность не подтверждается.

Локальный эффект продуктовых конкурентов почти везде отрицателен, от $-0.079$ до $0.003$ при
медиане $-0.029$, близкой к оценке МНК ($-0.035$). По Бонферрони он значим в 124 населенных пунктах, и все
значимые оценки отрицательны. Сильнее всего конкуренция проявляется на юге, в Ростовской (33
населенных пункта), Саратовской (17), Волгоградской (14), Воронежской (13) областях, а также на Севере
(рисунок \ref{fig:gwr_competitors}). Знак локального эффекта автомобильного трафика меняется по
территории (от $-0.074$ до $0.081$), значимы только положительные оценки в 35 населенных пунктах, из них 30 на
Севере (Мурманская и Архангельская области, Коми, Ямало-Ненецкий округ) и 5 в Нижегородской,
Вологодской и Кировской областях (рисунок \ref{fig:gwr_car_traffic}). Знак меняется и у ряда
других регрессоров, поэтому устойчивым по знаку можно считать прежде всего эффект конкурентов, а
также эффекты касс и доли больших магазинов, положительные во всех населенных пунктах.

\figuresingle{map_gwr_competitors.png}{Локальный эффект продуктовых конкурентов в GWR, значимые по Бонферрони}{fig:gwr_competitors}

\figuresingle{map_gwr_car_traffic.png}{Локальный эффект автомобильного трафика в GWR, значимые по Бонферрони}{fig:gwr_car_traffic}

\clearpage

% ============================================================
\section*{Заключение}

Выручка магазинов сети «Пятерочка» имеет выраженную и устойчивую во времени пространственную
структуру. Индекс Морана равен 0.58, кластеры высокой выручки образуют Московская и
Санкт-Петербургская агломерации с прилегающими областями, кластеры низкой - юг России и
Башкортостан; по двустороннему перестановочному тесту с поправкой FDR значимы локальные индексы
484 населенных пунктов из 933. Характеристики самих населенных пунктов объясняют лишь около половины этой зависимости.

Процедура выбора модели, сочетающая оба подхода, приводит к пространственной модели Дарбина с
матрицей экспоненциального затухания с масштабом 50 км, и ограничение $\theta = 0$ отвергается и
тестом отношения правдоподобия, и тестом Вальда, робастным к гетероскедастичности. Центральным
результатом является чувствительность к региональным факторам. С фиксированными эффектами
субъектов $\rho$ снижается с 0.47 до 0.20, а косвенные эффекты численности, автотрафика и
плотности сети становятся незначимыми. Отрицательная связь выручки с численностью соседних
населенных пунктов и положительная связь с их трафиком и плотностью сети наблюдаются в основной
спецификации, но не могут считаться установленными спилловерами между соседними рынками.
Пространственные панельные модели с двунаправленными фиксированными эффектами показывают, что
месячные отклонения выручки соседних населенных пунктов сильно связаны ($\rho \approx \lambda \approx 0.67$),
а динамическая модель с коррекцией смещения согласуется с одновременными пространственно
коррелированными шоками. Эффекты «субъект $\times$ месяц» снижают локальную автокорреляцию остатков
до значений не выше 0.08, хотя общая зависимость между населенными пунктами по CD-тесту сохраняется.
Пространственная зависимость и в уровнях, и в динамике согласуется с общими шоками и различиями
субъектов, но не доказывает их как единственный механизм.

Наиболее устойчивым результатом является локальная конкуренция. Рост среднего числа продуктовых
конкурентов в радиусе 500 м на единицу связан с выручкой магазина ниже на 2.2-3.3\% во всех 16
спецификациях. Это средняя условная связь на уровне населенного пункта, а не прогноз последствий
открытия одного конкурента. Свидетельств каннибализации между соседними магазинами сети
недостаточно, однако 95\%-интервал прямого эффекта плотности сети не исключает эффектов порядка
$\pm 2.5\%$. Разведочная географически взвешенная регрессия показывает, что эффект конкурентов
отрицателен почти во всех населенных пунктах и сильнее всего проявляется на юге России, но ее выводы о
значимости условны из-за пространственной зависимости остатков.

Результаты имеют ограничения:

\begin{enumerate}
    \item Изучаемая совокупность ограничена рынками, где сеть представлена не менее чем тремя
    магазинами;
    \item Размещение магазинов и число касс эндогенны, поэтому эффекты описывают связи, а не
    причинное влияние, а пространственный двухшаговый МНК инструментирует только $Wy$;
    \item Экзогенность матрицы весов является предположением, поскольку масштаб затухания выбран
    по правдоподобию;
    \item Цепной индекс цен не устраняет межрегиональных различий в уровне цен, их учитывают
    только эффекты субъектов;
    \item Стандартные ошибки пространственных панельных моделей не учитывают остаточную
    сериальную корреляцию;
    \item Оценки динамической панели обобщенным методом моментов не прошли $J$-тест, возможные
    причины - зависимость между населенными пунктами, пропущенная динамика и нарушение
    экзогенности региональных регрессоров.
\end{enumerate}

\clearpage

% ============================================================
\part*{Часть II. Прогнозирование розничного товарооборота с пространственными признаками}
\label{sec:forecast}

\section*{Постановка задачи и данные}

Цель эксперимента - проверить, могут ли пространственные признаки повысить точность прогноза розничного товарооборота (РТО) магазинов «Пятерочка».
В основной части исследования пространственные связи изучались эконометрическими методами.
Здесь использован более прикладной подход; характеристики соседних населенных пунктов добавлены к признакам моделей градиентного бустинга,
после чего сравнивается качество прогнозов.

В исходных данных 485\,082 наблюдения за январь 2023 - февраль 2025 года.
Каждая строка соответствует магазину в определенном месяце; задача - спрогнозировать РТО за март 2025 года.
Соревновательная оценка основана на средней абсолютной процентной ошибке (MAPE):
\begin{equation}
\operatorname{MAPE}=\frac{100}{n}\sum_{i=1}^n\left|\frac{\widehat y_i-y_i}{y_i}\right|,
\qquad
S=100\left(1-\frac{\min(\operatorname{MAPE},100)}{100}\right)^2.
\label{eq:score}
\end{equation}

В качестве отправной точки взят исходный ансамбль четырех моделей градиентного бустинга над решающими деревьями \cite{friedman2001}.
Он получил следующие результаты:

\begin{enumerate}
    \item 90,64 балла на публичной части тестового набора данных;
    \item 90,78 балла на итоговом тестовом наборе данных.
\end{enumerate}

Все результаты нового эксперимента далее приводятся для итогового тестового набора данных, если не указано иное. В таблице
ниже отдельно указаны результаты для обогащенного пространственными признаками ансамбля из градиентных бустингов и несколько иным видом
агрегации прогнозов базовых моделей.

\begin{table}[H]
\centering
\footnotesize
\caption{Сравнение исходного и итогового результатов}
\begin{tabularx}{\textwidth}{Xrr}\toprule
Вариант & Публичная часть & \shortstack{Итоговый тестовый\\набор данных}\\\midrule
Исходный ансамбль & 90,64 & 90,78\\
Ансамбль с пространственными признаками & - & \textbf{90,83}\\\bottomrule
\end{tabularx}
\tablesourcetext{расчеты авторов по данным X5 Group}
\end{table}

\subsection*{Исходное решение: отбор признаков и ансамблирование}

Процесс отбора признаков здесь достаточно непримечателен, хотя, в двух словах отметить стоит.
Для исходного признакового набора генерируется порядка 380 фичей (скользящие окна, лаги, агрегаты и прочие).
Далее проводится своего рода скорректированный под особенности задачи Forward Selection с нуля - важность считается по
значению функционала потерь (лосс-функции) базовой модели; например, CatBoost глубиной 5 в качестве базовой модели для всего признакового пространства,
а после итеративно отбирается по одной фиче из отсортированного по важности по вкладу в функционал потерь в пул выбранных фичей.
Каждые несколько фичей важность исходного набора (уже из оставшихся признаков) пересчитывается и процесс отбора сверху вниз запускается заново
(это что-то вроде страховки против дублирования сигналов, а также возможность учета перекрестных значимостей).
Само значение вклада в функционал потерь считается усреднением между 5 различными фолдами-месяцами в данных с различными весами.

Процесс подбора гиперпараметров аналогично проводился по взвешенному значению метрики на валидации,
но уже на 3-х фолдах (использовалась Optuna, метод - multivariate TPE) \cite{akiba2019, bergstra2011}.

Но более интересна здесь все же сама модель.

Следует отметить, что отбор признаков происходил для разных моделей вида градиентного бустинга;
под разными понимается обучение на разных таргетах - логарифм розничного товарооборота и лог-дифференцированное значение розничного товарооборота.
Соответственно, признаки были отобраны несколько разные для каждого из бустингов, но с некоторыми пересечениями.

После получения прогнозов на тестовых данных значения прогнозов усреднялись 4-мя различными методами усреднения, а после еще раз усреднялись
простым арифметическим средним для итогового предикта.

В общем, фактически, был реализован частный случай стекинга - блендинг с некоторой добавленной стохастичностью. Почему это работает?

Для простого арифметического среднего прогнозов $N$ моделей с одинаковой дисперсией ошибок
$\sigma^2$ и одинаковой попарной корреляцией ошибок $\rho_{\varepsilon}$ дисперсия ошибки
ансамбля равна:
\[
\operatorname{Var}\left(\frac{1}{N}\sum_{m=1}^{N}\varepsilon_m\right)
= \sigma^2\left(\rho_{\varepsilon} + \frac{1-\rho_{\varepsilon}}{N}\right),
\]
где $\varepsilon_m$ - ошибка прогноза $m$-й модели. При этих предпосылках меньшая корреляция
ошибок базовых моделей обеспечивает большее снижение дисперсии ошибки ансамбля.

В задаче, справедливости ради, прогнозы были достаточно схожими
(то есть, высокая корреляция перебивала полезный эффект от N моделей), поэтому эффект не был в достаточной мере большим,
но так или иначе среднее сглаживало перепрогнозы и недопрогнозы отдельных моделей.

Сама формула, скорее, схематично показывает причину улучшения итогового скора на лидерборде, так как,
с большой вероятностью, при ансамблировании ансамблей она становится заметно сложнее, но, тем не менее, по сути является таковой в обычном бэггинге.

Может ли быть оптимальное число моделей в таком ансамбле? Ответ на этот вопрос дает комбинаторная теория переобучения:

\begin{itemize}\item Если есть хорошая модель, а остальные - плохие, то хорошая модель будет превалировать и переобучения не будет;

\item Если есть много похожих друг на друга хороших моделей, переобучения тоже не будет, так как эффективная сложность
совокупности похожих друг на друга моделей невелика;

\item Но если получится так, что модели существенно различные и
все примерно одинаково плохие, то переобучение будет очень сильным, а сам эффект переобучения растет по мере роста числа моделей \cite{vorontsov2016}.
\end{itemize}

Отсюда следует, что, в целом, если все модели достаточно хорошие, имеет смысл увеличивать их количество; тем не менее,
на практике достаточно понять на валидации оптимальное минимальное число.

Схема алгоритма приведена на рисунке \ref{fig:forecast_scheme}. Так немного проще визуально понять, что именно и каким образом реализуется.

\figureforecast{forecast_model_scheme.jpg}{0.88}{Схема алгоритма исходного ансамбля}{fig:forecast_scheme}{составлено авторами}

\subsection*{Состав моделей и правило объединения прогнозов}

В новом эксперименте исходные признаки, преобразования целевой переменной и параметры обучения сохранены.
Это позволяет отдельно оценить эффект пространственных признаков, не смешивая его с результатом дополнительного подбора гиперпараметров.

\begin{table}[H]
\centering
\footnotesize
\caption{Модели исходного ансамбля}
\begin{tabularx}{\textwidth}{Xlr}\toprule
Модель & Целевая переменная & Число деревьев\\\midrule
CatBoost \cite{prokhorenkova2018} & $\log(1+y)$ & 6500\\
LightGBM (приросты) \cite{ke2017} & $\Delta\log y$ & 7000\\
LightGBM (уровни) & $\log(1+y)$ & 8000\\
XGBoost (приросты) \cite{chen2016} & $\Delta\log y$ & 3750\\\bottomrule
\end{tabularx}
\tablesourcetext{составлено авторами}
\end{table}

Прогнозы моделей, обученных на логарифме РТО, возвращаются к исходной шкале с помощью $\exp(\widehat z)-1$.
Для моделей, прогнозирующих логарифмический прирост, применяется $y_{i,t-1}\exp(\widehat g_{it})$, где $g_{it}=\log y_{it}-\log y_{i,t-1}$.
Ошибка обучения рассчитывается для преобразованной целевой переменной, а итоговая MAPE - для РТО в исходной шкале.

Для четырех прогнозов $p_1,\dots,p_4$ арифметическое и логарифмическое средние имеют вид:
\begin{equation}
A=\frac14\sum_{m=1}^4 p_m, \qquad
L=\exp\left(\frac14\sum_{m=1}^4\log(1+p_m)\right)-1.
\end{equation}

Дополнительно используются два пороговых правила, обозначенные $G$ и $H$. В $G$
арифметическое среднее заменяется медианой, когда $\max p_m/\min p_m>1{,}08$; в $H$ действует аналогичное правило,
но без модели LightGBM, прогнозирующей приросты. В принятой реализации итоговый прогноз равен:
\begin{equation}
P=\exp\left(\frac14\sum_{q\in\{A,L,G,H\}}\log(1+q)\right)-1.
\label{eq:meta}
\end{equation}
Именно это правило объединения прогнозов было сохранено при основном сравнении составов ансамбля.

\section*{Построение пространственных признаков}

Основная идея взята из пространственной модели SLX \cite{vega2015}, в которой наряду с собственными характеристиками объекта используются характеристики его соседей:
\begin{equation}
y=X\beta+WX\theta+\varepsilon.
\end{equation}
Здесь $W$ - матрица пространственных весов \cite{lesage2009}, а $WX$ - взвешенные характеристики соседних территорий. В отличие от линейной модели SLX,
в эксперименте $WX$ используется только как дополнительный набор входных признаков бустинга.
Прямые и косвенные причинные эффекты не оцениваются.

Прогноз строился для 18\,657 магазинов в 3\,360 населенных пунктах.
Пространственные признаки рассчитывались на уровне населенного пункта, а затем присоединялись к наблюдениям отдельных магазинов.
Таким образом, исходная панель магазинов не заменялась агрегированными данными.

Для построения $W$ проверены два варианта соседства: по 6 и по 12 ближайшим населенным пунктам (kNN6 и kNN12).
Расстояния между пунктами рассчитывались в метрической проекции Альберса с эллипсоидом Красовского.
Если один населенный пункт входит в список ближайших соседей другого, связь учитывается в обе стороны.
Затем веса нормируются так, чтобы их сумма для каждой строки равнялась единице:
\begin{equation}
w_{rv}=\frac{a_{rv}}{\sum_u a_{ru}}, \qquad w_{rr}=0,\qquad \sum_v w_{rv}=1.
\end{equation}

Для числового признака сначала находится среднее по магазинам населенного пункта, затем - взвешенное среднее по соседним пунктам:
\begin{equation}
\overline x_{rjt}=\frac{1}{|\mathcal I_{rj}^{\mathrm{obs}}(t)|}
\sum_{i\in\mathcal I_{rj}^{\mathrm{obs}}(t)}x_{ijt}, \qquad
(WX)_{rjt}=\sum_v w_{rv}\overline x_{vjt}.
\end{equation}
Магазины одного населенного пункта не считаются соседями друг для друга при расчете $WX$.
Для каждой модели формируется свой набор пространственных признаков на основе доступных исходных переменных;
идентификаторы и технические поля в него не включаются.
В итоге к CatBoost было добавлено 39 признаков, к LightGBM приростов - 22, к LightGBM уровней - 33, к XGBoost - 14.
Часть общих календарных признаков после пространственного усреднения остается неизменной и не дает нового сигнала.

Для прогноза марта использовались только сведения, доступные к концу февраля.
Фактический мартовский РТО не применялся при построении признаков; исторические лаги, напротив, были допустимы.
Это важно, чтобы избежать утечки целевой переменной. Пропущенные значения в пространственных признаках сохранялись без дополнительного заполнения.

\section*{Эксперименты и результаты}

\subsection*{Выбор пространственных признаков на исторических данных}

Для каждой из четырех моделей были проверены исходная версия без $WX$ и два варианта с kNN6 и kNN12.
Сначала варианты сравнивались по средней MAPE за август, ноябрь и декабрь 2024 года.
Затем выбранные пространственные версии и исходные модели проверялись на январе и феврале 2025 года.
В каждом случае обучение проводилось только на предшествующих месяцах.

\begin{table}[H]
\centering
\footnotesize
\caption{MAPE на этапе отбора пространственных признаков, \%}
\begin{tabularx}{\textwidth}{Xrrrr}\toprule
Модель & Без $WX$ & kNN6 & kNN12 & Выбор\\\midrule
CatBoost, уровни & 4,333 & 4,238 & 4,217 & $k=12$ \\
LGB, приросты & 5,335 & 4,179 & 4,080 & $k=12$ \\
LGB, уровни & 3,766 & 3,786 & 3,812 & $k=6$ \\
XGBoost, приросты & 4,864 & 4,744 & 4,664 & $k=12$ \\
\bottomrule\end{tabularx}
\tablesourcetext{расчеты авторов по данным X5 Group и GeoNames}
\end{table}

\begin{table}[H]
\centering
\footnotesize
\setlength{\tabcolsep}{4pt}
\caption{MAPE на проверке в январе и феврале 2025 года, \%}
\begin{tabularx}{\textwidth}{Xrrrrrr}\toprule
& & & \multicolumn{2}{c}{Январь} & \multicolumn{2}{c}{Февраль}\\
\cmidrule(lr){4-5}\cmidrule(lr){6-7}
Модель & Деревья & $k$ & До & С $WX$ & До & С $WX$\\\midrule
CatBoost, уровни & 6\,500 & 12 & 4,276 & 3,961 & 4,264 & 4,448 \\
LGB, приросты & 7\,000 & 12 & 4,003 & 4,000 & 5,122 & 4,689 \\
LGB, уровни & 8\,000 & 6 & 4,179 & 4,242 & 4,645 & 5,169 \\
XGBoost, приросты & 3\,750 & 12 & 4,030 & 4,064 & 4,428 & 4,714 \\
\bottomrule\end{tabularx}
\tablesourcetext{расчеты авторов по данным X5 Group и GeoNames}
\end{table}

Здесь результаты получились неоднозначными. CatBoost стал точнее в январе, но хуже в феврале.
LightGBM, прогнозирующий логарифмические приросты, улучшил MAPE в обоих месяцах; в феврале ошибка снизилась с 5,122 до 4,689\%.
Два других бустинга - LightGBM уровней и XGBoost - на этих месяцах уступили исходным версиям.
Поэтому по предварительному критерию для замены в ансамбле подходил только LightGBM приростов.

\subsection*{Результаты отдельных моделей на мартовском тестовом наборе данных}

\begin{table}[H]
\centering
\footnotesize
\caption{Итоговый балл отдельных моделей до и после добавления $WX$}
\begin{tabularx}{\textwidth}{Xrrr}\toprule
Модель & Без $WX$ & С $WX$ & Изменение\\\midrule
CatBoost, уровни & 90,55 & 90,39 & -0,16 \\
LGB, приросты & 90,32 & 90,24 & -0,08 \\
LGB, уровни & 90,57 & 90,55 & -0,02 \\
XGBoost, приросты & 90,45 & 90,58 & 0,13 \\
\bottomrule\end{tabularx}
\tablesourcetext{расчеты авторов по данным X5 Group и GeoNames}
\end{table}

На итоговом тестовом наборе данных за март результат улучшился только у XGBoost: с 90,45 до 90,58 балла.
CatBoost потерял 0,16 балла, LightGBM приростов - 0,08, LightGBM уровней - 0,02.
Следовательно, преимущество LightGBM приростов на январе и феврале не перенеслось на март, а XGBoost, наоборот, показал улучшение.

\figureforecast{forecast_model_comparison.png}{0.92}{Изменение ошибок на исторических месяцах и баллов на мартовском тестовом наборе данных}{fig:forecast_models}{составлено авторами по данным X5 Group и GeoNames}

Эти результаты наглядно показывают, что качество отдельных моделей на исторических месяцах не всегда соответствует их
результату на итоговом тестовом наборе данных в рамках соревнования. При этом важно оценивать не только отдельные прогнозы, но и их сочетание в ансамбле.

\subsection*{Сравнение составов ансамбля}

Чтобы сравнение было корректным, формула объединения прогнозов (\ref{eq:meta}) не менялась.
В каждом варианте заменялись только отдельные исходные модели на версии с пространственными признаками.

\begin{table}[H]
\centering
\footnotesize
\setlength{\tabcolsep}{4pt}
\caption{Сравнение составов ансамбля}
\begin{tabularx}{\textwidth}{Xrrr}\toprule
Модели с пространственными признаками & Январь, \% & Февраль, \% & Балл за март\\\midrule
Нет (исходный ансамбль) & 3,811945 & 4,265305 & 90,78\\
Только LightGBM приростов & 3,807809 & 4,209390 & 90,78\\
Все четыре модели & 3,763799 & 4,263002 & 90,79\\
Только XGBoost & 3,796333 & 4,291810 & 90,80\\
CatBoost и XGBoost & 3,754268 & 4,209174 & 90,76\\
XGBoost и LightGBM приростов & 3,789656 & 4,240780 & 90,82\\
XGBoost и LightGBM уровней & 3,805225 & 4,402870 & 90,81\\
XGBoost и оба LightGBM & 3,799520 & 4,349020 & \textbf{90,83}\\\bottomrule
\end{tabularx}
\tablesourcetext{расчеты авторов по данным X5 Group и GeoNames}
\end{table}

Предварительно выбранная замена LightGBM приростов уменьшила среднюю MAPE января и февраля с 4,038625 до 4,008600\%,
но на мартовском тестовом наборе данных балл остался равным 90,78. Сочетание CatBoost и XGBoost с $WX$, лучшее среди приведенных вариантов
по историческим месяцам, получило только 90,76.

Лучший мартовский результат, 90,83, получен при сохранении исходного CatBoost и замене XGBoost и двух LightGBM на версии с $WX$.
При этом средняя MAPE января и февраля у этого ансамбля составила 4,074270\% против 4,038625\% у исходного.
Таким образом, отбор по январю и февралю не привел бы к выбору именно этого состава.

Важно, что замена LightGBM, взятая отдельно, могла ухудшить его балл, но оказаться полезной в ансамбле.
По сути, это не является противоречием; вклад модели в итоговый прогноз зависит также от ошибок других компонентов.

\subsection*{Дополнительное сравнение способов усреднения}

\figureforecast{forecast_ensemble_grid.png}{0.92}{Результаты вариантов усреднения на мартовском тестовом наборе данных}{fig:forecast_ensembles}{составлено авторами по данным X5 Group и GeoNames}

Максимальные 90,83 балла получили два варианта:
исходное правило ансамблирования с $WX$ у XGBoost и обоих LightGBM (далее $A$) и вариант той же четверки
с пороговым выбором медианы при расхождении прогнозов более 1,08 (далее $B$). Простое арифметическое и логарифмическое
средние этих четырех моделей дали по 90,81 балла.

Также были рассмотрены два состава с результатом 90,82: $C$ - с $WX$ у XGBoost и
LightGBM приростов и пороговым правилом; $D$ - с $WX$ у XGBoost и LightGBM уровней, без LightGBM приростов.
Заключительная проверка касалась пяти способов объединения готовых ансамблей:

\begin{table}[H]
\centering
\footnotesize
\caption{Сравнение комбинаций готовых ансамблей}
\begin{tabularx}{\textwidth}{rXr}\toprule
№ & Способ объединения & Балл за март\\\midrule
1 & Арифметическое среднее $A,B$ & 90,83\\
2 & Логарифмическое среднее $A,B$ & 90,83\\
3 & Арифметическое среднее $A,B,C$ & 90,83\\
4 & Арифметическое среднее $A,B,D$ & 90,83\\
5 & Логарифмическое среднее $A,B,D$ & 90,83\\\bottomrule
\end{tabularx}
\tablesourcetext{расчеты авторов по данным X5 Group и GeoNames}
\end{table}

Ни один вариант не превысил результат 90,83 балла после округления. Для окончательного решения оставлен ансамбль $A$:
он обеспечивает тот же результат без дополнительного уровня усреднения.
Однако по опубликованным округленным оценкам нельзя утверждать, что его точная ошибка минимальна среди всех проверенных вариантов.

\section*{Пространственная зависимость ошибок и ограничения}

Для февраля 2025 года дополнительно рассчитан индекс Морана \cite{moran1950} по средним относительным ошибкам магазинов внутри населенных пунктов:
\begin{equation}
e_r=\operatorname{mean}_{i\in r}\left[\frac{y_i-\widehat y_i}{y_i}\right].
\end{equation}
Значимость проверялась перестановочным тестом с 999 перестановками. Здесь, в отличие от MAPE по магазинам, каждая территория имеет одинаковый вес.

\begin{table}[H]
\centering
\footnotesize
\caption{Индекс Морана для ошибок прогноза, февраль 2025 года}
\begin{tabularx}{\textwidth}{Xrrrr}\toprule
Модель & $k$ & До & С $WX$ & $p$-значение\\\midrule
CatBoost, уровни & 12 & 0,0657 & 0,0565 & 0,001 \\
LGB, приросты & 12 & 0,0559 & 0,0524 & 0,001 \\
LGB, уровни & 6 & 0,0960 & 0,0980 & 0,001 \\
XGBoost, приросты & 12 & 0,0710 & 0,0470 & 0,001 \\
\bottomrule\end{tabularx}
\tablesourcetext{расчеты авторов по данным X5 Group и GeoNames}
\end{table}

Пространственная зависимость ошибок сохраняется. Например, у XGBoost индекс Морана уменьшился с 0,0710 до 0,0470,
хотя февральская MAPE выросла с 4,428 до 4,714\%. Следовательно, уменьшение пространственной автокорреляции само по себе
еще не означает улучшения точности прогноза.

До основного эксперимента также рассматривались признаки на основе истории соседей, локальных пространственных статистик
и собственных векторов матрицы весов, а также варианты пространственных эконометрических моделей.
Некоторые из них улучшали показатели на отдельных исторических периодах, но не дали устойчивого выигрыша на мартовском тестовом наборе данных.
В отдельной проверке ранняя остановка при обучении также ухудшила качество относительно исходного режима.
Поэтому эти варианты не вошли в окончательный ансамбль.

Есть важное ограничение: после получения результатов на итоговом мартовском тестовом наборе данных они использовались при сравнении составов ансамбля.
Следовательно, 90,83 балла - результат на конкретном уже использованном тестовом наборе данных, а не независимая оценка качества на будущих месяцах.
Кроме того, без истинных значений мартовского РТО нельзя изучить корреляции ошибок базовых моделей и проверить статистическую
устойчивость полученного улучшения.

\section*{Выводы}

Добавление пространственных признаков позволило повысить результат исходного ансамбля с \textbf{90,78 до 90,83 балла}
на итоговом мартовском тестовом наборе данных. По формуле соревновательной метрики это приблизительно соответствует снижению MAPE с 4,7215 до 4,6952\%,
то есть на \textbf{0,0262 процентного пункта}, или примерно на 0,56\% относительно исходной ошибки.
Значения MAPE приближенные, поскольку баллы опубликованы с округлением.

При улучшении итогового результата с \textbf{90,78 до 90,83 балла} команда поднялась бы с \textbf{21-22-го на 12-е место}
в общем рейтинге и с \textbf{8-го на 5-е место} среди команд, приглашенных на защиту.
Таким образом, даже относительно небольшое улучшение качества прогноза позволило бы существенно повысить позицию команды в
итоговом рейтинге.

Наиболее полезной отдельной пространственной версией оказалась XGBoost. Но в лучшем ансамбле были также заменены
оба LightGBM, хотя по отдельности их баллы ухудшились. CatBoost остался без пространственных добавок. Эксперименты
с дополнительным усреднением не улучшили округленный итоговый результат.

Таким образом, пространственные характеристики соседних населенных пунктов могут давать дополнительный сигнал для прогнозирования РТО,
но эффект в этом эксперименте оказался небольшим и зависел от состава ансамбля; тем не менее, как таковой выигрыш от новых признаков заметен.

\clearpage

\begin{thebibliography}{99}
    \bibitem{vorontsov2016}
    Воронцов, К. В. Математические методы прогнозирования объемов продаж / К. В. Воронцов // ПостНаука : канал YouTube. - 2016. - 25 февраля. - URL: https://www.youtube.com/watch?v=u3lhVV3wAKk (дата обращения: 13.10.2026).

    \bibitem{demidova2021}
    Демидова, О. А. Пространственные аспекты оценки кривой заработной платы в России / О. А. Демидова, Е. А. Тимофеева // Журнал Новой экономической ассоциации. - 2021. - № 3 (51). - С. 69-101.

    \bibitem{akiba2019}
    Optuna: A Next-generation Hyperparameter Optimization Framework / T. Akiba, S. Sano, T. Yanase, T. Ohta, M. Koyama // Proceedings of the 25th ACM SIGKDD International Conference on Knowledge Discovery and Data Mining. - 2019. - P. 2623-2631.

    \bibitem{anselin1995}
    Anselin, L. Local Indicators of Spatial Association - LISA / L. Anselin // Geographical Analysis. - 1995. - Vol. 27, No. 2. - P. 93-115.

    \bibitem{anselin1996}
    Simple Diagnostic Tests for Spatial Dependence / L. Anselin, A. K. Bera, R. Florax, M. J. Yoon // Regional Science and Urban Economics. - 1996. - Vol. 26, No. 1. - P. 77-104.

    \bibitem{arellano1991}
    Arellano, M. Some Tests of Specification for Panel Data: Monte Carlo Evidence and an Application to Employment Equations / M. Arellano, S. Bond // Review of Economic Studies. - 1991. - Vol. 58, No. 2. - P. 277-297.

    \bibitem{benjamini1995}
    Benjamini, Y. Controlling the False Discovery Rate: A Practical and Powerful Approach to Multiple Testing / Y. Benjamini, Y. Hochberg // Journal of the Royal Statistical Society. Series B. - 1995. - Vol. 57, No. 1. - P. 289-300.

    \bibitem{benjamini2001}
    Benjamini, Y. The Control of the False Discovery Rate in Multiple Testing under Dependency / Y. Benjamini, D. Yekutieli // Annals of Statistics. - 2001. - Vol. 29, No. 4. - P. 1165-1188.

    \bibitem{bergstra2011}
    Algorithms for Hyper-Parameter Optimization / J. Bergstra, R. Bardenet, Y. Bengio, B. Kegl // Advances in Neural Information Processing Systems. - 2011. - Vol. 24. - P. 2546-2554.

    \bibitem{blundell1998}
    Blundell, R. Initial Conditions and Moment Restrictions in Dynamic Panel Data Models / R. Blundell, S. Bond // Journal of Econometrics. - 1998. - Vol. 87, No. 1. - P. 115-143.

    \bibitem{chen2016}
    Chen, T. XGBoost: A Scalable Tree Boosting System / T. Chen, C. Guestrin // Proceedings of the 22nd ACM SIGKDD International Conference on Knowledge Discovery and Data Mining. - 2016. - P. 785-794.

    \bibitem{driscoll1998}
    Driscoll, J. C. Consistent Covariance Matrix Estimation with Spatially Dependent Panel Data / J. C. Driscoll, A. C. Kraay // Review of Economics and Statistics. - 1998. - Vol. 80, No. 4. - P. 549-560.

    \bibitem{elhorst2014}
    Elhorst, J. P. Spatial Econometrics: From Cross-Sectional Data to Spatial Panels / J. P. Elhorst. - Berlin : Springer, 2014.

    \bibitem{fotheringham2002}
    Fotheringham, A. S. Geographically Weighted Regression: The Analysis of Spatially Varying Relationships / A. S. Fotheringham, C. Brunsdon, M. Charlton. - Chichester : Wiley, 2002.

    \bibitem{friedman2001}
    Friedman, J. H. Greedy Function Approximation: A Gradient Boosting Machine / J. H. Friedman // Annals of Statistics. - 2001. - Vol. 29, No. 5. - P. 1189-1232.

    \bibitem{getis1992}
    Getis, A. The Analysis of Spatial Association by Use of Distance Statistics / A. Getis, J. K. Ord // Geographical Analysis. - 1992. - Vol. 24, No. 3. - P. 189-206.

    \bibitem{vega2015}
    Halleck Vega, S. The SLX Model / S. Halleck Vega, J. P. Elhorst // Journal of Regional Science. - 2015. - Vol. 55, No. 3. - P. 339-363.

    \bibitem{holmes2011}
    Holmes, T. J. The Diffusion of Wal-Mart and Economies of Density / T. J. Holmes // Econometrica. - 2011. - Vol. 79, No. 1. - P. 253-302.

    \bibitem{huff1964}
    Huff, D. L. Defining and Estimating a Trading Area / D. L. Huff // Journal of Marketing. - 1964. - Vol. 28, No. 3. - P. 34-38.

    \bibitem{jia2008}
    Jia, P. What Happens When Wal-Mart Comes to Town: An Empirical Analysis of the Discount Retailing Industry / P. Jia // Econometrica. - 2008. - Vol. 76, No. 6. - P. 1263-1316.

    \bibitem{ke2017}
    LightGBM: A Highly Efficient Gradient Boosting Decision Tree / G. Ke, Q. Meng, T. Finley, T. Wang, W. Chen, W. Ma, Q. Ye, T.-Y. Liu // Advances in Neural Information Processing Systems. - 2017. - Vol. 30. - P. 3146-3154.

    \bibitem{kelejian1998}
    Kelejian, H. H. A Generalized Spatial Two-Stage Least Squares Procedure for Estimating a Spatial Autoregressive Model with Autoregressive Disturbances / H. H. Kelejian, I. R. Prucha // Journal of Real Estate Finance and Economics. - 1998. - Vol. 17, No. 1. - P. 99-121.

    \bibitem{kelejian1999}
    Kelejian, H. H. A Generalized Moments Estimator for the Autoregressive Parameter in a Spatial Model / H. H. Kelejian, I. R. Prucha // International Economic Review. - 1999. - Vol. 40, No. 2. - P. 509-533.

    \bibitem{kelejian2010}
    Kelejian, H. H. Specification and Estimation of Spatial Autoregressive Models with Autoregressive and Heteroskedastic Disturbances / H. H. Kelejian, I. R. Prucha // Journal of Econometrics. - 2010. - Vol. 157, No. 1. - P. 53-67.

    \bibitem{lee2010dynamic}
    Lee, L. F. A Spatial Dynamic Panel Data Model with Both Time and Individual Fixed Effects / L. F. Lee, J. Yu // Econometric Theory. - 2010. - Vol. 26, No. 2. - P. 564-597.

    \bibitem{lee2010}
    Lee, L. F. Estimation of Spatial Autoregressive Panel Data Models with Fixed Effects / L. F. Lee, J. Yu // Journal of Econometrics. - 2010. - Vol. 154, No. 2. - P. 165-185.

    \bibitem{lesage2009}
    LeSage, J. Introduction to Spatial Econometrics / J. LeSage, R. K. Pace. - Boca Raton : Chapman and Hall/CRC, 2009.

    \bibitem{lesage2014}
    LeSage, J. P. Spatial Econometric Panel Data Model Specification: A Bayesian Approach / J. P. LeSage // Spatial Statistics. - 2014. - Vol. 9. - P. 122-145.

    \bibitem{leung2000}
    Leung, Y. Statistical Tests for Spatial Nonstationarity Based on the Geographically Weighted Regression Model / Y. Leung, C.-L. Mei, W.-X. Zhang // Environment and Planning A. - 2000. - Vol. 32, No. 1. - P. 9-32.

    \bibitem{moran1950}
    Moran, P. A. P. Notes on Continuous Stochastic Phenomena / P. A. P. Moran // Biometrika. - 1950. - Vol. 37, No. 1/2. - P. 17-23.

    \bibitem{nickell1981}
    Nickell, S. Biases in Dynamic Models with Fixed Effects / S. Nickell // Econometrica. - 1981. - Vol. 49, No. 6. - P. 1417-1426.

    \bibitem{pesaran2021}
    Pesaran, M. H. General Diagnostic Tests for Cross-Sectional Dependence in Panels / M. H. Pesaran // Empirical Economics. - 2021. - Vol. 60, No. 1. - P. 13-50.

    \bibitem{prokhorenkova2018}
    CatBoost: Unbiased Boosting with Categorical Features / L. Prokhorenkova, G. Gusev, A. Vorobev, A. V. Dorogush, A. Gulin // Advances in Neural Information Processing Systems. - 2018. - Vol. 31. - P. 6638-6648.

    \bibitem{roodman2009}
    Roodman, D. A Note on the Theme of Too Many Instruments / D. Roodman // Oxford Bulletin of Economics and Statistics. - 2009. - Vol. 71, No. 1. - P. 135-158.

    \bibitem{yu2008}
    Yu, J. Quasi-Maximum Likelihood Estimators for Spatial Dynamic Panel Data with Fixed Effects When Both n and T Are Large / J. Yu, R. de Jong, L. F. Lee // Journal of Econometrics. - 2008. - Vol. 146, No. 1. - P. 118-134.
\end{thebibliography}

\end{document}
